\documentclass[10pt, letterpaper]{article}

% Inhaltsverzeichnis für Pakettypen (nur für Übersicht im Header, wird nicht im Dokument angezeigt)
% 1. Seitenlayout und Ränder
% 2. Sprache und Zeichensatz
% 3. Mathematik und Theorem-Umgebungen
% 4. Eigene Makros
% 5. Diagramme und Grafiken
% 6. Tabellen und Aufzählungen
% 7. Inhaltsverzeichnis
% 8. Abschnittsüberschriften
% 9. Abstrakt-Umgebung
% 10. Todos/Notizen
% 11. Rahmen/Box-Umgebungen
% 12. Python-Integration
% 13. Literaturverwaltung
% 14. Hyperlinks
% 15. Absatzeinstellungen
% 16. Umgebungen
% 17. Titel und Autor

% --- 0. Allgemeine Pakete ---
\usepackage{booktabs}
\usepackage{makecell}
\usepackage{slashed}

% --- 1. Seitenlayout und Ränder ---
\usepackage[margin=3cm]{geometry}

% --- 2. Sprache und Zeichensatz ---
\usepackage[english]{babel}
\usepackage[T1]{fontenc}
\usepackage[utf8]{inputenc}

% --- 3. Mathematik und Theorem-Umgebungen ---
\usepackage{amsmath, amssymb, amsthm}
\usepackage{mathrsfs}
\DeclareMathOperator{\WF}{WF}

% --- 4. Eigene Makros ---
\usepackage{xcolor}
\newcommand{\SKP}{\langle\cdot,\cdot\rangle}
\newcommand{\R}{\mathbb{R}}
\newcommand{\N}{\mathbb{N}}
\newcommand{\Q}{\mathbb{Q}}
\newcommand{\Z}{\mathbb{Z}}
\newcommand{\C}{\mathbb{C}}
\newcommand{\entwurf}[1]{\textcolor{red}{#1}}
\newcommand{\GL}{\mathrm{GL}}
\newcommand{\SL}{\mathrm{SL}}
\newcommand{\SO}{\mathrm{SO}}
\newcommand{\SU}{\mathrm{SU}}
\newcommand{\Sp}{\mathrm{Sp}}
\newcommand{\Ogroup}{\mathrm{O}} % \O ist schon für skandinavisches Ø reserviert
\newcommand{\U}{\mathrm{U}}

% --- 5. Diagramme und Grafiken ---
\usepackage{graphicx}
\graphicspath{{../../Images/}}
\usepackage[export]{adjustbox}
\usepackage{tikz}
\usetikzlibrary{decorations.pathreplacing, arrows.meta, positioning}
\usepackage{tikz-cd}

% --- 6. Tabellen und Aufzählungen ---
\usepackage{enumitem}
\setlist[itemize]{left=0.5cm}

\newenvironment{romanenum}[1][]
  {%
    \ifx&#1&
    \else
      \textbf{#1}\quad
    \fi
    \begin{enumerate}[label=\roman*)]
  }
  {%
    \end{enumerate}%
  }

% --- 7. Inhaltsverzeichnis ---
\usepackage{tocloft}
\renewcommand{\cftsecfont}{\footnotesize}
\renewcommand{\cftsubsecfont}{\footnotesize}
\renewcommand{\cftsubsubsecfont}{\footnotesize}
\renewcommand{\cftsecpagefont}{\footnotesize}
\renewcommand{\cftsubsecpagefont}{\footnotesize}
\renewcommand{\cftsubsubsecpagefont}{\footnotesize}
\usepackage{etoc}

% --- 8. Abschnittsüberschriften ---
\usepackage{titlesec}
\titleformat{\section}{\normalfont\large\bfseries}{\thesection}{1em}{}
\titleformat{\subsection}{\normalfont\normalsize\bfseries}{\thesubsection}{0.5em}{}
\titleformat{\subsubsection}{\normalfont\normalsize\bfseries}{\thesubsubsection}{0.5em}{}
\setcounter{secnumdepth}{4}

% --- 9. Abstrakt-Umgebung ---
\usepackage{changepage}
\renewenvironment{abstract}
  {
    \begin{adjustwidth}{1.5cm}{1.5cm}
    \small
    \textsc{Abstract. –}%
  }
  {
    \end{adjustwidth}
  }

% --- 10. Todos/Notizen ---
\usepackage{todonotes}
% Hellblau definieren
\definecolor{hellblau}{rgb}{0.3,0.6,1}

% Eigene Umgebung »notiz«
\newenvironment{note}{\par\begingroup\color{hellblau}\itshape}{\par\endgroup}

% --- 11. Rahmen/Box-Umgebungen ---
\usepackage{mdframed}
\usepackage{tcolorbox}
\colorlet{shadecolor}{gray!25}

\newenvironment{customTheorem}
  {\vspace{10pt}%
   \begin{mdframed}[
     backgroundcolor=gray!20,
     linewidth=0pt,
     innertopmargin=10pt,
     innerbottommargin=10pt,
     skipabove=\dimexpr\topsep+\ht\strutbox\relax,
     skipbelow=\topsep,
   ]}
  {\end{mdframed}
   \vspace{10pt}%
  }

% --- 12. Python-Integration ---
% (Deaktiviert in dieser Version, aktiviere bei Bedarf)
% \usepackage{pythontex}
% \usepackage[makestderr]{pythontex}

% --- 13. Literaturverwaltung ---
\usepackage{csquotes}
\usepackage[backend=biber, style=alphabetic, citestyle=alphabetic]{biblatex}
\addbibresource{bibliography.bib}

% --- 14. Hyperlinks ---
\usepackage{hyperref}
\hypersetup{
  colorlinks   = true,
  urlcolor     = blue,
  linkcolor    = blue,
  citecolor    = blue,
  frenchlinks  = true
}

% --- 15. Absatzeinstellungen ---
\usepackage[parfill]{parskip}
\sloppy

% --- 16. Umgebungen ---
\usepackage{thmtools}

\newcommand{\CustomHeading}[3]{%
  \par\medskip\noindent%
  \textbf{#1 #2} \textnormal{(#3)}.\enskip%
}

\newenvironment{DEF}[2]{\CustomHeading{Definition}{#1}{#2}}{}
\newenvironment{PROP}[2]{\CustomHeading{Proposition}{#1}{#2}}{}
\newenvironment{THEO}[2]{\CustomHeading{Theorem}{#1}{#2}}{}
\newenvironment{LEM}[2]{\CustomHeading{Lemma}{#1}{#2}}{}
\newenvironment{KORO}[2]{\CustomHeading{Corollar}{#1}{#2}}{}
\newenvironment{REM}[2]{\CustomHeading{Remark}{#1}{#2}}{}
\newenvironment{EXA}[2]{\CustomHeading{Example}{#1}{#2}}{}
\newenvironment{STUD}[2]{\CustomHeading{Study}{#1}{#2}}{}
\newenvironment{CONC}[2]{\CustomHeading{Concept}{#1}{#2}}{}
\newenvironment{OTH}[2]{\CustomHeading{Other}{#1}{#2}}{}
\newenvironment{EXE}[2]{\CustomHeading{Exercise}{#1}{#2}}{}
\newenvironment{MOT}[2]{\CustomHeading{Motivation}{#1}{#2}}{}
\newenvironment{PROOF}[2]{\CustomHeading{Proof}{#1}{#2}}{}
\newenvironment{STR}[2]{\CustomHeading{Structure}{#1}{#2}}{}


% --- Unit Umgebung ---
\usepackage{mdframed}
\newmdenv[
  linewidth=1pt,
  topline=false,
  bottomline=false,
  rightline=false,
  leftmargin=0cm,
  rightmargin=0cm,
  skipabove=10pt,
  skipbelow=10pt,
  innertopmargin=0.5\baselineskip,
  innerbottommargin=0.5\baselineskip,
  backgroundcolor=gray!10,
  linecolor=gray
]{unitbox}

\newenvironment{unit}[1]
  {\begin{unitbox}\textbf{Unit #1}\par\smallskip}
  {\end{unitbox}}


% --- 17. Titel und Autor ---
\title{Mein Titel}
\author{Tim Jaschik}
\date{\today}

\begin{document}

\maketitle
\rule{\textwidth}{0.5pt}
\begin{abstract}
Kurze Beschreibung …
\end{abstract}
\rule{\textwidth}{0.5pt}
\vspace{0.5cm}

\tableofcontents

\pagebreak

\section{Definitionen}

\subsection{Matrix Lie Gruppen}

\begin{DEF}{LIE-H15-01-01}{Allgemeine Lineare Gruppen}
Für $n\in\mathbb{N}$ sei
\[
\GL(n;\R)=\{A\in M_n(\R)\mid \det A\neq 0\},\qquad
\GL(n;\C)=\{A\in M_n(\C)\mid \det A\neq 0\}.
\]
\end{DEF}

\begin{DEF}{LIE-H15-01-02}{Matrixraum und Konvergenz}
$M_n(\C)$ ist der Raum aller $n\times n$-Matrizen (identifiziert mit $\C^{n^2}$).
Eine Folge $(A_m)\subset M_n(\C)$ konvergiert gegen $A\in M_n(\C)$, wenn für alle $1\le j,k\le n$ die Einträge $(A_m)_{jk}\to A_{jk}$ konvergieren.
\end{DEF}

\begin{DEF}{LIE-H15-01-03}{Matrix-Lie-Gruppe}
Eine Matrix-Lie-Gruppe ist eine Untergruppe $G\subset \GL(n;\C)$ mit:
Konvergiert $(A_m)\subset G$ gegen $A\in M_n(\C)$, dann gilt: Entweder $A\in G$ oder $A$ ist nicht invertierbar.
Äquivalent: $G$ ist eine abgeschlossene Untergruppe von $\GL(n;\C)$ (in der Teilraumtopologie).
\end{DEF}

\begin{DEF}{LIE-H15-01-04}{Lineare und spezielle lineare Gruppen}
$\GL(n;\C)$ und $\GL(n;\R)$ sind Matrix-Lie-Gruppen.
Die speziellen linearen Gruppen
\[
\SL(n;\C)=\{A\in M_n(\C)\mid \det A=1\},\qquad
\SL(n;\R)=\{A\in M_n(\R)\mid \det A=1\}
\]
sind wegen Stetigkeit der Determinante abgeschlossen, also Matrix-Lie-Gruppen.
\end{DEF}

\begin{DEF}{LIE-H15-01-05}{Unitäre, spezielle unitäre und orthogonale Gruppen}
$A\in M_n(\C)$ heißt unitär, wenn $A^*A=I$ (also $A^{-1}=A^*$) oder in Komponenten gilt:
$$
\sum_{l=1}^{n} \overline{A_{l j}} A_{l k}=\delta_{j k}
$$
Dann
\[
\U(n)=\{A\in M_n(\C)\mid A^*A=I\},\qquad
\SU(n)=\{A\in \U(n)\mid \det A=1\}.
\]
Für reelle Matrizen ist $A$ orthogonal, wenn $A^{\mathrm{tr}}A=I$; damit
\[
\Ogroup(n)=\{A\in M_n(\R)\mid A^{\mathrm{tr}}A=I\},\qquad
\SO(n)=\{A\in \Ogroup (n)\mid \det A=1\}.
\]
\end{DEF}

\begin{DEF}{LIE-H15-01-06}{Komplex orthogonale Gruppen}
Mit der bilinearen Form $(x,y)=\sum_j x_jy_j$ gilt
\[
\Ogroup(n;\C)=\{A\in M_n(\C)\mid A^{\mathrm{tr}}A=I\},\qquad
\SO(n;\C)=\{A\in \Ogroup (n;\C)\mid \det A=1\}.
\]
Dies sind Matrix-Lie-Gruppen (nicht zu verwechseln mit $\U(n)$).
\end{DEF}

\begin{DEF}{LIE-H15-01-07}{Verallgemeinerte orthogonale (Lorentz-)Gruppen}
Für $n,k\in\N$ und $x,y\in\R^{n+k}$:
\[
[x,y]_{n,k}=x_1y_1+\cdots+x_ny_n - x_{n+1}y_{n+1}-\cdots - x_{n+k}y_{n+k}.
\]
Dann
\[
\Ogroup(n;k)=\{A\in \GL(n{+}k;\R)\mid [Ax,Ay]_{n,k}=[x,y]_{n,k}\ \forall x,y\},\quad
\SO(n;k)=\{A\in \Ogroup (n;k)\mid \det A=1\}.
\]
Mit
\[
g=\mathrm{diag}(\underbrace{1,\dots,1}_{n},\underbrace{-1,\dots,-1}_{k})
\]
gilt äquivalent: $A\in\Ogroup(n;k)\iff A^{\mathrm{tr}}gA=g$. Der Fall $\Ogroup(3;1)$ ist die Lorentzgruppe.
\end{DEF}

\begin{DEF}{LIE-H15-01-08}{Symplektische Gruppen (reell, komplex, kompakt)}
Auf $\R^{2n}$:
\[
\omega(x,y)=\sum_{j=1}^n\big(x_j y_{n+j} - x_{n+j} y_j\big),\qquad
\Omega=\begin{pmatrix}0&I\\ -I&0\end{pmatrix}.
\]
Reell:
\[
\Sp(n;\R)=\{A\in \GL(2n;\R)\mid \omega(Ax,Ay)=\omega(x,y)\ \forall x,y\}.
\]
Äquivalent:
\[
A\in\Sp(n;\R)\iff -\Omega A^{\mathrm{tr}}\Omega = A^{-1}.
\]
Komplex:
\[
\Sp(n;\C)=\{A\in \GL(2n;\C)\mid -\Omega A^{\mathrm{tr}}\Omega = A^{-1}\}.
\]
Kompakt:
\[
\Sp(n)=\Sp(n;\C)\cap \U(2n).
\]
\end{DEF}

\begin{DEF}{LIE-H15-01-09}{Euklidische und Poincaré-Gruppe}
Euklidisch: Elemente sind $\{x,R\}$ mit $x\in\R^n$, $R\in\Ogroup(n)$ und Wirkung $\{x,R\}\cdot y=Ry+x$.
Gruppenverknüpfung und Inverses:
\[
\{x_1,R_1\}\{x_2,R_2\}=\{x_1+R_1x_2,\ R_1R_2\},\qquad
\{x,R\}^{-1}=\{-R^{-1}x,\ R^{-1}\}.
\]
Matrixdarstellung (Isomorphie zu Untergruppe von $\GL(n{+}1;\R)$):
\[
\begin{pmatrix}
R & x\\
0 & 1
\end{pmatrix},\quad R\in\Ogroup (n),\ x\in\R^n.
\]
Poincaré (inhomogene Lorentzgruppe): $\R^{n+1}\rtimes \Ogroup (n;1)$ mit Darstellungen
\[
\begin{pmatrix}
A & x\\
0 & 1
\end{pmatrix},\quad A\in\Ogroup (n;1),\ x\in\R^{n+1}.
\]
\end{DEF}

\begin{DEF}{LIE-H15-01-10}{Heisenberg-Gruppe}
\[
H=\left\{
\begin{pmatrix}
1 & a & b\\
0 & 1 & c\\
0 & 0 & 1
\end{pmatrix}
\ \middle|\ a,b,c\in\R
\right\}\subset \GL(3;\R),\qquad
\begin{pmatrix}
1 & a & b\\
0 & 1 & c\\
0 & 0 & 1
\end{pmatrix}^{-1}
=
\begin{pmatrix}
1 & -a & ac-b\\
0 & 1 & -c\\
0 & 0 & 1
\end{pmatrix}.
\]

$\R^*,\ \C^*,\ S^1,\ \R,\ \R^n$ als Matrixgruppen]
\[
\R^* \cong \GL(1;\R),\qquad \C^* \cong \GL(1;\C),\qquad S^1=\{z\in\C:\ |z|=1\}\cong \U(1).
\]
\[
(\R,+)\cong \GL(1;\R)^+,\quad x\mapsto [e^x],\qquad
(\R^n,+)\cong\left\{\mathrm{diag}(e^{x_1},\dots,e^{x_n})\mid x\in\R^n\right\}.
\]
\end{DEF}

\begin{DEF}{LIE-H15-01-11}{Kompakte symplektische Gruppe als „unitär über $\mathbb{H}$“}
Auf $\C^{2n}$ sei der konjugiert-lineare Operator
\[
J(\alpha,\beta)=(-\overline{\beta},\,\overline{\alpha}),\qquad
\omega(z,w)=\langle Jz,w\rangle,\qquad J^2=-I.
\]
Dann
\[
\Sp(n)=\{U\in \U(2n)\mid UJ=JU\},
\]
d.\,h.\ $\Sp(n)$ ist genau die Menge der unitären Operatoren, die mit $J$ kommutieren; dies erlaubt die Interpretation von $\Sp(n)$ als „unitäre Gruppe über den Quaternionen“.
\end{DEF}

\begin{DEF}{LIE-H15-01-34}{Konjugierte Lineare Abbildung $J$ zur Charakterisierung von $\Sp(n)$}
Of the groups introduced in the preceding subsections, the compact symplectic group $\operatorname{Sp}(n):=\operatorname{Sp}(n ; \mathbb{C}) \cap \mathrm{U}(2 n)$ is the most mysterious. In this section, we attempt to understand the structure of $\operatorname{Sp}(n)$ and to show that it can be understood as being the "unitary group over the quaternions."

Since the definition of $\operatorname{Sp}(n)$ involves unitarity, it is convenient to express the bilinear form $\omega$ on $\mathbb{C}^{2 n}$ in terms of the inner product $\langle\cdot, \cdot\rangle$, rather than in terms of the bilinear form ( $\cdot, \cdot$ ), as we did in Sect. 1.2.4. To this end, define a conjugate-linear map $J: \mathbb{C}^{2 n} \rightarrow \mathbb{C}^{2 n}$ by
$$
J(\alpha, \beta)=(-\bar{\beta}, \bar{\alpha})
$$
where $\alpha$ and $\beta$ are in $\mathbb{C}^{n}$ and $(\alpha, \beta)$ is in $\mathbb{C}^{2 n}$. We can easily check that for all $z, w \in \mathbb{C}^{2 n}$, we have
$$
\omega(z, w)=\langle J z, w\rangle
$$
Recall that we take our inner product to be conjugate linear in the first factor; since $J$ is also conjugate linear, $\langle J z, w\rangle$ is actually linear in $z$. We may easily check that
$$
\langle J z, w\rangle=-\overline{\langle z, J w\rangle}=-\langle J w, z\rangle
$$
for all $z, w \in \mathbb{C}^{2 n}$ and that
$$
J^{2}=-I .
$$
\end{DEF}

\begin{DEF}{LIE-H15-01-45}{Kompakte Matrix Lie Gruppen}
A matrix Lie group $G \subset \mathrm{GL}(n ; \mathbb{C})$ is said to be compact if it is compact in the usual topological sense as a subset of $M_{n}(\mathbb{C}) \cong \mathbb{R}^{2 n^{2}}$.
\end{DEF}

\begin{DEF}{LIE-H15-01-46}{Charakterisierung von Kompakte Matrix Lie Gruppen}
In light of the Heine-Borel theorem (Theorem 2.41 in [Rud1]), a matrix Lie group $G$ is compact if and only if it is closed (as a subset of $M_{n}(\mathbb{C})$, not just as a subset of $\mathrm{GL}(n ; \mathbb{C})$ ) and bounded. Explicitly, this means that $G$ is compact if and only if (1) whenever $A_{m} \in G$ and $A_{m} \rightarrow A$, then $A$ is in $G$, and (2) there exists a constant $C$ such that for all $A \in G$, we have $\left|A_{j k}\right| \leq C$ for all $1 \leq j, k \leq n$.
\end{DEF}

\begin{DEF}{LIE-H15-01-49}{Wegzusammenhängende Matrix Lie Gruppen}
A matrix Lie group $G$ is said to be connected if for all $A$ and $B$ in $G$, there exists a continuous path $A(t), a \leq t \leq b$, lying in $G$ with $A(a)=A$ and $A(b)=B$. For any matrix Lie group $G$, the identity component of $G$, denoted $G_{0}$, is the set of $A \in G$ for which there exists a continuous path $A(t), a \leq t \leq b$, lying in $G$ with $A(a)=I$ and $A(b)=A$.
\end{DEF}

\begin{DEF}{LIE-H15-01-51}{Charakterisierung von WZSH Matrix Lie Gruppen}
To show that a matrix Lie group $G$ is connected, it suffices to show that each $A \in G$ can be connected to the identity by a continuous path lying in $G$.
\end{DEF}

\begin{DEF}{LIE-H15-01-64}{Einfach Zusammenhängende Matrix Lie Gruppe}
A matrix Lie group $G$ is said to be simply connected if it is connected and, in addition, every loop in $G$ can be shrunk continuously to a point in $G$.
\end{DEF}

\begin{DEF}{LIE-H15-01-65}{Charakterisierung von EZSH Matrix Lie Gruppen}
More precisely, assume that $G$ is connected. Then $G$ is simply connected if for every continuous path $A(t), 0 \leq t \leq 1$, lying in $G$ and with $A(0)=A(1)$, there exists a continuous function $A(s, t), 0 \leq s, t \leq 1$, taking values in $G$ and having the following properties: (1) $A(s, 0)=A(s, 1)$ for all $s$, (2) $A(0, t)=A(t)$, and (3) $A(1, t)=A(1,0)$ for all $t$.

One should think of $A(t)$ as a loop and $A(s, t)$ as a family of loops, parameterized by the variable $s$ which shrinks $A(t)$ to a point. Condition 1 says that for each value of the parameter $s$, we have a loop; Condition 2 says that when $s=0$ the loop is the specified loop $A(t)$; and Condition 3 says that when $s=1$ our loop is a point. The condition of simple connectedness is important because for simply connected groups, there is a particularly close relationship between the group and the Lie algebra.
\end{DEF}

\begin{DEF}{LIE-H15-01-68}{Reell-Projektiver Raum}
The real projective space of dimension $n$, denoted $\mathbb{R} P^{n}$, is the set of lines through the origin in $\mathbb{R}^{n+1}$. Since each line through the origin intersects the unit sphere exactly twice, we may think of $\mathbb{R} P^{n}$ as the unit sphere $S^{n}$ with "antipodal" points $u$ and $-u$ identified. Using the second description, we think of points in $\mathbb{R} P^{n}$ as pairs $\{u,-u\}$, with $u \in S^{n}$.
\end{DEF}

\begin{DEF}{LIE-H15-01-71}{Charakterisierung von $\R P^n$ durch obere Halbsphäre}
Let $D^{n}$ denote the closed upper hemisphere in $S^{n}$, that is, the set of points $u \in S^{n}$ with $u_{n+1} \geq 0$. Then $\pi$ maps $D^{n}$ onto $\mathbb{R} P^{n}$, since at least one of $u$ and $-u$ is in $D^{n}$. The restriction of $\pi$ to $D^{n}$ is injective except on the equator, that is, the set of $u \in S^{n}$ with $u_{n+1}=0$. If $u$ is in the equator, then $-u$ is also in the equator, and $\pi(-u)=\pi(u)$. Thus, we may also think of $\mathbb{R} P^{n}$ as the upper hemisphere $D^{n}$, with antipodal points on the equator identified (Figure 1.2).

\includegraphics[scale=0.2, center]{2025_08_28_fc86fd4ade69f41b5842g-033}

Fig. 1.2 The space $\mathbb{R} P^{n}$ is the upper hemisphere with antipodal points on the equator identified. The indicated path from $u$ to $-u$ corresponds to a loop in $\mathbb{R} P^{n}$ that cannot be shrunk to a point.
\end{DEF}

\begin{DEF}{LIE-H15-01-72}{Charakterisierung von $\R P^n$ durch Scheibe mit antipodal identifizierten Rand}
We may now make one last identification using the projection $P$ of $\mathbb{R}^{n+1}$ onto $\mathbb{R}^{n}$. (That is to say, $P$ is the map sending $\left(x_{1}, \ldots, x_{n}, x_{n+1}\right)$ to $\left(x_{1}, \ldots, x_{n}\right)$.) The restriction of $P$ to $D^{n}$ is a continuous bijection between $D^{n}$ and the closed unit ball $B^{n}$ in $\mathbb{R}^{n}$, with the equator in $D^{n}$ mapping to the boundary of the ball. Thus, our last model of $\mathbb{R} P^{n}$ is the closed unit ball $B^{n} \subset \mathbb{R}^{n}$, with antipodal points on the boundary of $B^{n}$ identified.
\end{DEF}

\begin{DEF}{LIE-H15-01-75}{Lie Gruppen Homomorphismus}
Let $G$ and $H$ be matrix Lie groups. A map $\Phi$ from $G$ to $H$ is called a Lie group homomorphism if (1) $\Phi$ is a group homomorphism and (2) $\Phi$ is continuous. If, in addition, $\Phi$ is one-to-one and onto and the inverse map $\Phi^{-1}$ is continuous, then $\Phi$ is called a Lie group isomorphism.
\end{DEF}

\begin{DEF}{LIE-H15-01-77}{Lie Gruppen Isomorphismus}
Note that the inverse of a Lie group isomorphism is continuous (by definition) and a group homomorphism (by elementary group theory), and thus a Lie group isomorphism. If $G$ and $H$ are matrix Lie groups and there exists a Lie group isomorphism from $G$ to $H$, then $G$ and $H$ are said to be isomorphic, and we write $G \cong H$.
\end{DEF}

\begin{DEF}{LIE-H15-01-88}{Lie Gruppen}
A Lie group is a smooth manifold $G$ which is also a group and such that the group product
$$
G \times G \rightarrow G
$$
and the inverse map $G \rightarrow G$ are smooth.
\end{DEF}

\subsection{Exponentiale}

\begin{DEF}{LIE-H15-02-01}{Exponentialabbildung für Matrizen}
If $X$ is an $n \times n$ matrix, we define the exponential of $X$, denoted $e^{X}$ or $\exp X$, by the usual power series
$$
e^{X}=\sum_{m=0}^{\infty} \frac{X^{m}}{m!},
$$
where $X^{0}$ is defined to be the identity matrix $I$ and where $X^{m}$ is the repeated matrix product of $X$ with itself.
\end{DEF}

\begin{DEF}{LIE-H15-02-03}{Hilbert-Schmidt-Norm für Matrizen}
For any $X \in M_{n}(\mathbb{C})$, we define
$$
\|X\|=\left(\sum_{j, k=1}^{n}\left|X_{j k}\right|^{2}\right)^{1 / 2}
$$
The quantity $\|X\|$ is called the Hilbert-Schmidt norm of $X$.

The Hilbert-Schmidt norm may be computed in a basis-independent way as
$$
\|X\|=\left(\operatorname{trace}\left(X^{*} X\right)\right)^{1 / 2}
$$
\end{DEF}

\begin{DEF}{LIE-H15-02-07}{Charakterisierung von Folgen Konvergenz via Hilbert-Schmidt Norm}
If $X_{m}$ is a sequence of matrices, then it is easy to see that $X_{m}$ converges to a matrix $X$ in the sense of Definition 1.3 if and only if $\left\|X_{m}-X\right\| \rightarrow 0$ as $m \rightarrow \infty$.
\end{DEF}

\begin{DEF}{LIE-H15-02-20}{Logarithmus für Matrizen}
For an $n \times n$ matrix $A$, define $\log A$ by
$$
\log A=\sum_{m=1}^{\infty}(-1)^{m+1} \frac{(A-I)^{m}}{m}
$$
whenever the series converges.
\end{DEF}

\begin{DEF}{LIE-H15-02-33}{1-Parameter Untergruppen von $\GL(n;\C)$}
A function $A: \mathbb{R} \rightarrow \mathrm{GL}(n ; \mathbb{C})$ is called a one-parameter subgroup of $\mathrm{GL}(n ; \mathbb{C})$ if
\begin{enumerate}
  \item $A$ is continuous,
  \item $A(0)=I$,
  \item $A(t+s)=A(t) A(s)$ for all $t, s \in \mathbb{R}$.
\end{enumerate}
\end{DEF}

\begin{DEF}{LIE-H15-02-41}{Positive selbstadjungierte Matrix}
If $P$ is a self-adjoint $n \times n$ matrix (i.e., $P^{*}=P$ ), we say that $P$ is positive if $\langle v, P v\rangle>0$ for all nonzero $v \in \mathbb{C}^{n}$. It is easy to check that a self-adjoint matrix $P$ is positive if and only if all the eigenvalues of $P$ are positive.
\end{DEF}

\subsection{Lie Algebren}

\begin{DEF}{LIE-H15-03-01}{Lie Algebra}
A finite-dimensional real or complex Lie algebra is a finitedimensional real or complex vector space $\mathfrak{g}$, together with a map $[\cdot, \cdot]$ from $\mathfrak{g} \times \mathfrak{g}$ into $\mathfrak{g}$, with the following properties:
\begin{enumerate}
  \item $[\cdot, \cdot]$ is bilinear.
  \item $[\cdot, \cdot]$ is skew symmetric: $[X, Y]=-[Y, X]$ for all $X, Y \in \mathfrak{g}$.
  \item The Jacobi identity holds:
  $$
  [X,[Y, Z]]+[Y,[Z, X]]+[Z,[X, Y]]=0
  $$
  for all $X, Y, Z \in \mathfrak{g}$.
\end{enumerate}
Two elements $X$ and $Y$ of a Lie algebra $\mathfrak{g}$ commute if $[X, Y]=0$. A Lie algebra $\mathfrak{g}$ is commutative if $[X, Y]=0$ for all $X, Y \in \mathfrak{g}$.

The map $[\cdot, \cdot]$ is referred to as the bracket operation on $\mathfrak{g}$. Note also that Condition 2 implies that $[X, X]=0$ for all $X \in \mathfrak{g}$. The bracket operation on a Lie algebra is not, in general associative; nevertheless, the Jacobi identity can be viewed as a substitute for associativity.
\end{DEF}

\begin{DEF}{LIE-H15-03-06}{Subalgebren von Lie Algebren (Reell, Komplex)}
A subalgebra of a real or complex Lie algebra $\mathfrak{g}$ is a subspace $\mathfrak{h}$ of $\mathfrak{g}$ such that $\left[H_{1}, H_{2}\right] \in \mathfrak{h}$ for all $H_{1}$ and $H_{2} \in \mathfrak{h}$. If $\mathfrak{g}$ is a complex Lie algebra and $\mathfrak{h}$ is a real subspace of $\mathfrak{g}$ which is closed under brackets, then $\mathfrak{h}$ is said to be a real subalgebra of $\mathfrak{g}$.
\end{DEF}

\begin{DEF}{LIE-H15-03-07}{Ideale von Lie Algebren}
A subalgebra $\mathfrak{h}$ of a Lie algebra $\mathfrak{g}$ is said to be an ideal in $\mathfrak{g}$ if $[X, H] \in \mathfrak{h}$ for all $X$ in $\mathfrak{g}$ and $H$ in $\mathfrak{h}$.
\end{DEF}

\begin{DEF}{LIE-H15-03-08}{Zentrum einer Lie Algebra}
The center of a Lie algebra $\mathfrak{g}$ is the set of all $X \in \mathfrak{g}$ for which $[X, Y]=0$ for all $Y \in \mathfrak{g}$.
\end{DEF}

\begin{DEF}{LIE-H15-03-09}{Lie Algebra Homomorphismus}
If $\mathfrak{g}$ and $\mathfrak{h}$ are Lie algebras, then a linear map $\phi: \mathfrak{g} \rightarrow \mathfrak{h}$ is called a Lie algebra homomorphism if $\phi([X, Y])=[\phi(X), \phi(Y)]$ for all $X, Y \in \mathfrak{g}$. If, in addition, $\phi$ is one-to-one and onto, then $\phi$ is called a Lie algebra isomorphism. A Lie algebra isomorphism of a Lie algebra with itself is called a Lie algebra automorphism.
\end{DEF}

\begin{DEF}{LIE-H15-03-10}{Adjungierte Abbildung von Lie Algebren}
If $\mathfrak{g}$ is a Lie algebra and $X$ is an element of $\mathfrak{g}$, define a linear map $\operatorname{ad}_{X}: \mathfrak{g} \rightarrow \mathfrak{g}$ by
$$
\operatorname{ad}_{X}(Y)=[X, Y] .
$$
The map $X \mapsto \operatorname{ad}_{X}$ is the adjoint map or adjoint representation.
\end{DEF}

\begin{DEF}{LIE-H15-03-14}{Direkte Summe von Lie Algebren}
If $\mathfrak{g}_{1}$ and $\mathfrak{g}_{2}$ are Lie algebras, the direct sum of $\mathfrak{g}_{1}$ and $\mathfrak{g}_{2}$ is the vector space direct sum of $\mathfrak{g}_{1}$ and $\mathfrak{g}_{2}$, with bracket given by
$$
\left[\left(X_{1}, X_{2}\right),\left(Y_{1}, Y_{2}\right)\right]=\left(\left[X_{1}, Y_{1}\right],\left[X_{2}, Y_{2}\right]\right) .
$$
\end{DEF}

\begin{DEF}{LIE-H15-03-15}{Lie Algebra als direkte Summe von Unteralgebren}
If $\mathfrak{g}$ is a Lie algebra and $\mathfrak{g}_{1}$ and $\mathfrak{g}_{2}$ are subalgebras, we say that $\mathfrak{g}$ decomposes as the Lie algebra direct sum of $\mathfrak{g}_{1}$ and $\mathfrak{g}_{2}$ if $\mathfrak{g}$ is the direct sum of $\mathfrak{g}_{1}$ and $\mathfrak{g}_{2}$ as vector spaces and $\left[X_{1}, X_{2}\right]=0$ for all $X_{1} \in \mathfrak{g}_{1}$ and $X_{2} \in \mathfrak{g}_{2}$.
\end{DEF}

\begin{DEF}{LIE-H15-03-17}{Strukturkonstanten von endlichen Lie Algebren}
Let $\mathfrak{g}$ be a finite-dimensional real or complex Lie algebra, and let $X_{1}, \ldots, X_{N}$ be a basis for $\mathfrak{g}$ (as a vector space). Then the unique constants $c_{j k l}$ such that
$$
\left[X_{j}, X_{k}\right]=\sum_{l=1}^{N} c_{j k l} X_{l}
$$
are called the structure constants of $\mathfrak{g}$ (with respect to the chosen basis).
\end{DEF}

\begin{DEF}{LIE-H15-03-19}{Irreduzible Lie Algebra}
A Lie algebra $\mathfrak{g}$ is called irreducible if the only ideals in $\mathfrak{g}$ are $\mathfrak{g}$ and $\{0\}$.
\end{DEF}

\begin{DEF}{LIE-H15-03-20}{Einfache Lie Algebra}
A Lie algebra $\mathfrak{g}$ is called simple if it is irreducible and $\operatorname{dim} \mathfrak{g} \geq 2$.
\end{DEF}

\begin{DEF}{LIE-H15-03-24}{Kommutator-Ideal in Lie Algebren}
If $\mathfrak{g}$ is a Lie algebra, then the commutator ideal in $\mathfrak{g}$, denoted $[\mathfrak{g}, \mathfrak{g}]$, is the space of linear combinations of commutators, that is, the space of elements $Z$ in $\mathfrak{g}$ that can be expressed as
$$
Z=c_{1}\left[X_{1}, Y_{1}\right]+\cdots+c_{m}\left[X_{m}, Y_{m}\right]
$$
for some constants $c_{j}$ and vectors $X_{j}, Y_{j} \in \mathfrak{g}$.

For any $X$ and $Y$ in $\mathfrak{g}$, the commutator $[X, Y]$ is in $[\mathfrak{g}, \mathfrak{g}]$. This holds, in particular, if $X$ is in $[\mathfrak{g}, \mathfrak{g}]$, showing that $[\mathfrak{g}, \mathfrak{g}]$ is an ideal in $\mathfrak{g}$.
\end{DEF}

\begin{DEF}{LIE-H15-03-25}{Abgeleitete Reihe von Lie Algebren}
For any Lie algebra $\mathfrak{g}$, we define a sequence of subalgebras $\mathfrak{g}_{0}, \mathfrak{g}_{1}, \mathfrak{g}_{2}, \ldots$ of $\mathfrak{g}$ inductively as follows: $\mathfrak{g}_{0}=\mathfrak{g}, \mathfrak{g}_{1}=\left[\mathfrak{g}_{0}, \mathfrak{g}_{0}\right], \mathfrak{g}_{2}=\left[\mathfrak{g}_{1}, \mathfrak{g}_{1}\right]$, etc. These subalgebras are called the derived series of $\mathfrak{g}$. A Lie algebra $\mathfrak{g}$ is called solvable if $\mathfrak{g}_{j}=\{0\}$ for some $j$.
\end{DEF}

\begin{DEF}{LIE-H15-03-26}{Aufgelöste Lie Algebra}
For any Lie algebra $\mathfrak{g}$, we define a sequence of subalgebras $\mathfrak{g}_{0}, \mathfrak{g}_{1}, \mathfrak{g}_{2}, \ldots$ of $\mathfrak{g}$ inductively as follows: $\mathfrak{g}_{0}=\mathfrak{g}, \mathfrak{g}_{1}=\left[\mathfrak{g}_{0}, \mathfrak{g}_{0}\right], \mathfrak{g}_{2}=\left[\mathfrak{g}_{1}, \mathfrak{g}_{1}\right]$, etc. These subalgebras are called the derived series of $\mathfrak{g}$. A Lie algebra $\mathfrak{g}$ is called solvable if $\mathfrak{g}_{j}=\{0\}$ for some $j$.
\end{DEF}

\begin{DEF}{LIE-H15-03-28}{Obere zentrale Reihe einer Lie Algebra}
For any Lie algebra $\mathfrak{g}$, we define a sequence of ideals $\mathfrak{g}^{j}$ in $\mathfrak{g}$ inductively as follows. We set $\mathfrak{g}^{0}=\mathfrak{g}$ and then define $\mathfrak{g}^{j+1}$ to be the space of linear combinations of commutators of the form $[X, Y]$ with $X \in \mathfrak{g}$ and $Y \in \mathfrak{g}^{j}$. These algebras are called the upper central series of $\mathfrak{g}$. A Lie algebra $\mathfrak{g}$ is said to be nilpotent if $\mathfrak{g}^{j}=\{0\}$ for some $j$.
Equivalently, $\mathfrak{g}^{j}$ is the space spanned by all $j$ th-order commutators,
$$
\left[X_{1},\left[X_{2},\left[X_{3}, \ldots\left[X_{j}, X_{j+1}\right] \ldots\right]\right]\right] .
$$
\end{DEF}

\begin{DEF}{LIE-H15-03-29}{Nilpotente Lie Algebren}
For any Lie algebra $\mathfrak{g}$, we define a sequence of ideals $\mathfrak{g}^{j}$ in $\mathfrak{g}$ inductively as follows. We set $\mathfrak{g}^{0}=\mathfrak{g}$ and then define $\mathfrak{g}^{j+1}$ to be the space of linear combinations of commutators of the form $[X, Y]$ with $X \in \mathfrak{g}$ and $Y \in \mathfrak{g}^{j}$. These algebras are called the upper central series of $\mathfrak{g}$. A Lie algebra $\mathfrak{g}$ is said to be nilpotent if $\mathfrak{g}^{j}=\{0\}$ for some $j$.
\end{DEF}

\begin{DEF}{LIE-H15-03-33}{Lie Algebra einer Matrix Lie Gruppe}
Let $G$ be a matrix Lie group. The Lie algebra of $G$, denoted $\mathfrak{g}$, is the set of all matrices $X$ such that $e^{t X}$ is in $G$ for all real numbers $t$.

Equivalently, $X$ is in $\mathfrak{g}$ if and only if the entire one-parameter subgroup generated by $X$ lies in $G$.
\end{DEF}

\begin{DEF}{LIE-H15-03-40}{Komplexe Matrix Lie Gruppe}
A matrix Lie group $G$ is said to be complex if its Lie algebra $\mathfrak{g}$ is a complex subspace of $M_{n}(\mathbb{C})$, that is, if $i X \in \mathfrak{g}$ for all $X \in \mathfrak{g}$.
\end{DEF}

\begin{DEF}{LIE-H15-03-65}{Adjunigerte Abbildung auf Lie Algebra von MLG}
Let $G$ be a matrix Lie group, with Lie algebra $\mathfrak{g}$. Then for each $A \in G$, define a linear map $\operatorname{Ad}_{A}: \mathfrak{g} \rightarrow \mathfrak{g}$ by the formula
$$
\operatorname{Ad}_{A}(X)=A X A^{-1}
$$
\end{DEF}

\begin{DEF}{LIE-H15-03-66}{Adjunigerte Abbildung auf Lie Algebra von MLG}
Let $G$ be a matrix Lie group, with Lie algebra $\mathfrak{g}$. Then for each $A \in G$, define a linear map $\operatorname{Ad}_{A}: \mathfrak{g} \rightarrow \mathfrak{g}$ by the formula
$$
\operatorname{Ad}_{A}(X)=A X A^{-1}
$$
\end{DEF}

\begin{DEF}{LIE-H15-03-69}{Induzierte lineare Abbildung von Lie Algebra in Menge der Endomorphismen der Lie Algebra}
By Theorem 3.28, there is an associated real linear map $X \rightarrow \mathrm{ad}_{X}$ from the Lie algebra of $G$ to the Lie algebra of $\mathrm{GL}(\mathfrak{g})$ (i.e., from $\mathfrak{g}$ to $\mathrm{gl}(\mathfrak{g})$ ), with the property that
$$
e^{\mathrm{ad}_{X}}=\operatorname{Ad}_{e^{X}} .
$$
Here, $\operatorname{gl}(\mathfrak{g})$ is the Lie algebra of $\operatorname{GL}(\mathfrak{g})$, namely the space of all linear maps of $\mathfrak{g}$ to itself.
\end{DEF}

\begin{DEF}{LIE-H15-03-74}{Komplexifikation von end-dim reellen Vektorräumen}
If $V$ is a finite-dimensional real vector space, then the complexification of $V$, denoted $V_{\mathbb{C}}$, is the space of formal linear combinations
$$
v_{1}+i v_{2}
$$
with $v_{1}, v_{2} \in V$. This becomes a real vector space in the obvious way and becomes a complex vector space if we define
$$
i\left(v_{1}+i v_{2}\right)=-v_{2}+i v_{1} .
$$

We could more pedantically define $V_{\mathbb{C}}$ to be the space of ordered pairs ( $v_{1}, v_{2}$ ) with $v_{1}, v_{2} \in V$, but this is notationally cumbersome. It is straightforward to verify that the above definition really makes $V_{\mathbb{C}}$ into a complex vector space. We will regard $V$ as a real subspace of $V_{\mathbb{C}}$ in the obvious way.
\end{DEF}

\begin{DEF}{LIE-H15-03-83}{Exponential Abbildung von Lie Algebra in Matrix Lie Gruppen}
If $G$ is a matrix Lie group with Lie algebra $\mathfrak{g}$, then the exponential map for $G$ is the map
$$
\exp : \mathfrak{g} \rightarrow G .
$$
That is to say, the exponential map for $G$ is the matrix exponential restricted to the Lie algebra $\mathfrak{g}$ of $G$. 

We have shown (Theorem 2.10) that every matrix in $\mathrm{GL}(n ; \mathbb{C})$ is the exponential of some $n \times n$ matrix. Nevertheless, if $G \subset \operatorname{GL}(n ; \mathbb{C})$ is a closed subgroup, there may exist $A$ in $G$ such that there is no $X$ in the Lie algebra $\mathfrak{g}$ of $G$ with $\exp X=A$.
\end{DEF}

\section{Exampel}

\subsection{Matrix Lie Gruppen}

\begin{EXA}{LIE-H15-01-16}{Menge der invertierbaren Matrizen mit rationalen Einträgen keine MLG}
An example of a subgroup of $\mathrm{GL}(n ; \mathbb{C})$ which is not closed (and hence is not a matrix Lie group) is the set of all $n \times n$ invertible matrices with rational entries. This set is, in fact, a subgroup of $\mathrm{GL}(n ; \mathbb{C})$, but not a closed subgroup. That is, one can (easily) have a sequence of invertible matrices with rational entries converging to an invertible matrix with some irrational entries. (In fact, every real invertible matrix is the limit of some sequence of invertible matrices with rational entries.)
\end{EXA}

\begin{EXA}{LIE-H15-01-17}{Irrationale Linie in einem Torus keine MLG}
Another example of a group of matrices which is not a matrix Lie group is the following subgroup of $\operatorname{GL}(2 ; \mathbb{C})$. Let $a$ be an irrational real number and let
$$
G=\left\{\left.\left(\begin{array}{cc}
e^{i t} & 0 \\
0 & e^{i t a}
\end{array}\right) \right\rvert\, t \in \mathbb{R}\right\} .
$$
Clearly, $G$ is a subgroup of $\mathrm{GL}(2 ; \mathbb{C})$. According to Exercise 10, the closure of $G$ is the group
$$
\bar{G}=\left\{\left.\left(\begin{array}{cc}
e^{i \theta} & 0 \\
0 & e^{i \phi}
\end{array}\right) \right\rvert\, \theta, \phi \in \mathbb{R}\right\} .
$$
The group $G$ inside $\bar{G}$ is known as an "irrational line in a torus":


\includegraphics[scale=0.2, center]{2025_08_28_fc86fd4ade69f41b5842g-018}

Fig. 1.1 A small portion of the group $G$ inside $\bar{G}$ (left) and a larger portion (right)
\end{EXA}

\begin{EXA}{LIE-H15-01-18}{$\GL(n;\C)$ ist MLG}
The general linear groups (over $\mathbb{R}$ or $\mathbb{C}$ ) are themselves matrix Lie groups. Of course, $\mathrm{GL}(n ; \mathbb{C})$ is a subgroup of itself. Furthermore, if $A_{m}$ is a sequence of matrices in $\mathrm{GL}(n ; \mathbb{C})$ and $A_{m}$ converges to $A$, then by the definition of $\mathrm{GL}(n ; \mathbb{C})$, either $A$ is in $\mathrm{GL}(n ; \mathbb{C})$, or $A$ is not invertible.
\end{EXA}

\begin{EXA}{LIE-H15-01-19}{$\GL(n;\R)$ ist MLG}
Moreover, $\mathrm{GL}(n ; \mathbb{R})$ is a subgroup of $\mathrm{GL}(n ; \mathbb{C})$, and if $A_{m} \in \mathrm{GL}(n ; \mathbb{R})$ and $A_{m}$ converges to $A$, then the entries of $A$ are real. Thus, either $A$ is not invertible or $A \in \mathrm{GL}(n ; \mathbb{R})$.
\end{EXA}

\begin{EXA}{LIE-H15-01-20}{$\SL(n;\R,\C)$ sind MLG}
The special linear group (over $\mathbb{R}$ or $\mathbb{C}$ ) is the group of $n \times n$ invertible matrices (with real or complex entries) having determinant one. Both of these are subgroups of $\mathrm{GL}(n ; \mathbb{C})$. Furthermore, if $A_{n}$ is a sequence of matrices with determinant one and $A_{n}$ converges to $A$, then $A$ also has determinant one, because the determinant is a continuous function. Thus, $\mathrm{SL}(n ; \mathbb{R})$ and $\mathrm{SL}(n ; \mathbb{C})$ are matrix Lie groups.
\end{EXA}

\begin{EXA}{LIE-H15-01-21}{$\U(n), \SU(n)$ sind MLG}
An $n \times n$ complex matrix $A$ is said to be unitary if the column vectors of $A$ are orthonormal, that is, if
$$
\sum_{l=1}^{n} \overline{A_{l j}} A_{l k}=\delta_{j k}
$$
We may rewrite (1.2) as
$$
\sum_{l=1}^{n}\left(A^{*}\right)_{j l} A_{l k}=\delta_{j k}
$$
where $\delta_{j k}$ is the Kronecker delta, equal to 1 if $j=k$ and equal to zero if $j \neq k$. Here $A^{*}$ is the adjoint of $A$, defined by
$$
\left(A^{*}\right)_{j k}=\overline{A_{k j}}
$$
Equation (1.3) says that $A^{*} A=I$; thus, we see that $A$ is unitary if and only if $A^{*}=A^{-1}$. In particular, every unitary matrix is invertible.

The adjoint operation on matrices satisfies $(A B)^{*}=B^{*} A^{*}$. From this, we can see that if $A$ and $B$ are unitary, then
$$
(A B)^{*}(A B)=B^{*} A^{*} A B=B^{-1} A^{-1} A B=I
$$
showing that $A B$ is also unitary. Furthermore, since $\left(A A^{-1}\right)^{*}=I^{*}=I$, we see that $\left(A^{-1}\right)^{*} A^{*}=I$, which shows that $\left(A^{-1}\right)^{*}=\left(A^{*}\right)^{-1}$. Thus, if $A$ is unitary, we have
$$
\left(A^{-1}\right)^{*} A^{-1}=\left(A^{*}\right)^{-1} A^{-1}=\left(A A^{*}\right)^{-1}=I,
$$
showing that $A^{-1}$ is again unitary.

Thus, the collection of unitary matrices is a subgroup of $\operatorname{GL}(n ; \mathbb{C})$. We call this group the unitary group and we denote it by $\mathrm{U}(n)$. 


We may also define the special unitary group $\mathrm{SU}(n)$, the subgroup of $\mathrm{U}(n)$ consisting of unitary matrices with determinant 1 . It is easy to check that both $\mathrm{U}(n)$ and $\mathrm{SU}(n)$ are closed subgroups of $\mathrm{GL}(n ; \mathbb{C})$ and thus matrix Lie groups.
\end{EXA}

\begin{EXA}{LIE-H15-01-22}{$\Ogroup(n), \SO(n)$ sind MLG}
Meanwhile, let $\langle\cdot, \cdot\rangle$ denote the standard inner product on $\mathbb{C}^{n}$, given by
$$
\langle x, y\rangle=\sum_{j} \overline{x_{j}} y_{j}
$$
(Note that we put the conjugate on the first factor in the inner product.) By Proposition A.8, we have
$$
\langle x, A y\rangle=\left\langle A^{*} x, y\right\rangle
$$
for all $x, y \in \mathbb{C}^{n}$. Thus,
$$
\langle A x, A y\rangle=\left\langle A^{*} A x, y\right\rangle,
$$
from which we can see that if $A$ is unitary, then $A$ preserves the inner product on $\mathbb{C}^{n}$, that is,
$$
\langle A x, A y\rangle=\langle x, y\rangle
$$
for all $x$ and $y$. Conversely, if $A$ preserves the inner product, we must have $\left\langle A^{*} A x, y\right\rangle=\langle x, y\rangle$ for all $x, y$. It is not hard to see that this condition holds only if $A^{*} A=I$. Thus, an equivalent characterization of unitarity is that $A$ is unitary if and only if $A$ preserves the standard inner product on $\mathbb{C}^{n}$.

Finally, for any matrix $A$, we have that $\operatorname{det} A^{*}=\overline{\operatorname{det} A}$. Thus, if $A$ is unitary, we have
$$
\operatorname{det}\left(A^{*} A\right)=|\operatorname{det} A|^{2}=\operatorname{det} I=1
$$
Hence, for all unitary matrices $A$, we have $|\operatorname{det} A|=1$.

In a similar fashion, an $n \times n$ real matrix $A$ is said to be orthogonal if the column vectors of $A$ are orthonormal. As in the unitary case, we may give equivalent versions of this condition. The only difference is that if $A$ is real, $A^{*}$ is the same as the transpose $A^{t r}$ of $A$, given by
$$
\left(A^{t r}\right)_{j k}=A_{k j} .
$$
Thus, $A$ is orthogonal if and only if $A^{t r}=A^{-1}$, and this holds if and only if $A$ preserves the inner product on $\mathbb{R}^{n}$. Since $\operatorname{det}\left(A^{t r}\right)=\operatorname{det} A$, if $A$ is orthogonal, we have
$$
\operatorname{det}\left(A^{t r} A\right)=\operatorname{det}(A)^{2}=\operatorname{det}(I)=1,
$$
so that $\operatorname{det}(A)= \pm 1$. The collection of all orthogonal matrices forms a closed subgroup of $\mathrm{GL}(n ; \mathbb{C})$, which we call the orthogonal group and denote by $\mathrm{O}(n)$. 

The set of $n \times n$ orthogonal matrices with determinant one is the special orthogonal group, denoted $\mathrm{SO}(n)$. Geometrically, elements of $\mathrm{SO}(n)$ are rotations, while the elements of $\mathrm{O}(n)$ are either rotations or combinations of rotations and reflections.
\end{EXA}

\begin{EXA}{LIE-H15-01-23}{$\Ogroup(n,\C), \SO(n,\C)$ sind MLG}
Consider now the bilinear form ( $\cdot, \cdot$ ) on $\mathbb{C}^{n}$ defined by
$$
(x, y)=\sum_{j} x_{j} y_{j} .
$$
This form is not an inner product (Sect. A.6) because, for example, it is symmetric rather than conjugate-symmetric. The set of all $n \times n$ complex matrices $A$ which preserve this form (i.e., such that $(A x, A y)=(x, y)$ for all $x, y \in \mathbb{C}^{n}$ ) is the complex orthogonal group $\mathrm{O}(n ; \mathbb{C})$, and it is a subgroup of $\mathrm{GL}(n ; \mathbb{C})$. Since there are no conjugates in the definition of the form ( $\cdot, \cdot$ ), we have
$$
(x, A y)=\left(A^{t r} x, y\right),
$$
for all $x, y \in \mathbb{C}^{n}$, where on the right-hand side of the above relation, we have $A^{t r}$ rather than $A^{*}$. Repeating the arguments for the case of $\mathrm{O}(n)$, but now allowing complex entries in our matrices, we find that an $n \times n$ complex matrix $A$ is in $\mathrm{O}(n ; \mathbb{C})$ if and only if $A^{t r} A=I$, that $\mathrm{O}(n ; \mathbb{C})$ is a matrix Lie group, and that $\operatorname{det} A= \pm 1$ for all $A$ in $\mathrm{O}(n ; \mathbb{C})$. Note that $\mathrm{O}(n ; \mathbb{C})$ is not the same as the unitary group $\mathrm{U}(n)$. 

The group $\mathrm{SO}(n ; \mathbb{C})$ is defined to be the set of all $A$ in $\mathrm{O}(n ; \mathbb{C})$ with $\operatorname{det} A=1$ and it is also a matrix Lie group.
\end{EXA}

\begin{EXA}{LIE-H15-01-24}{$\Ogroup(n,k)$ ist MLG und Lorentz-Gruppe}
Let $n$ and $k$ be positive integers, and consider $\mathbb{R}^{n+k}$. Define a symmetric bilinear form $[\cdot, \cdot]_{n, k}$ on $\mathbb{R}^{n+k}$ by the formula

$$
[x, y]_{n, k}=x_{1} y_{1}+\cdots+x_{n} y_{n}-x_{n+1} y_{n+1}-\cdots-x_{n+k} y_{n+k}
$$

The set of $(n+k) \times(n+k)$ real matrices $A$ which preserve this form (i.e., such that $[A x, A y]_{n, k}=[x, y]_{n, k}$ for all $x, y \in \mathbb{R}^{n+k}$ ) is the generalized orthogonal group $\mathrm{O}(n ; k)$. It is a subgroup of $\mathrm{GL}(n+k ; \mathbb{R})$ and a matrix Lie group.


Of particular interest in physics is the Lorentz group $\mathrm{O}(3 ; 1)$. We also define $\mathrm{SO}(n ; k)$ to be the subgroup of $\mathrm{O}(n ; k)$ consisting of elements with determinant 1 .
\end{EXA}

\begin{EXA}{LIE-H15-01-27}{$\Sp(2n,\R)$ ist MLG}
Consider the skew-symmetric bilinear form $B$ on $\mathbb{R}^{2 n}$ defined as follows:
$$
\omega(x, y)=\sum_{j=1}^{n}\left(x_{j} y_{n+j}-x_{n+j} y_{j}\right) .
$$
The set of all $2 n \times 2 n$ matrices $A$ which preserve $\omega$ (i.e., such that $\omega(A x, A y)= \omega(x, y)$ for all $x, y \in \mathbb{R}^{2 n}$ ) is the real symplectic group $\operatorname{Sp}(n ; \mathbb{R})$, and it is a closed subgroup of $\operatorname{GL}(2 n ; \mathbb{R})$. (Some authors refer to the group we have just defined as $\operatorname{Sp}(2 n ; \mathbb{R})$ rather than $\operatorname{Sp}(n ; \mathbb{R})$.) If $\Omega$ is the $2 n \times 2 n$ matrix
$$
\Omega=\left(\begin{array}{cc}
0 & I \\
-I & 0
\end{array}\right),
$$
then
$$
\omega(x, y)=\langle x, \Omega y\rangle .
$$
From this, it is not hard to show that a $2 n \times 2 n$ real matrix $A$ belongs to $\operatorname{Sp}(n ; \mathbb{R})$ if and only if
$$
-\Omega A^{t r} \Omega=A^{-1} .
$$
Taking the determinant of this identity gives $\operatorname{det} A=(\operatorname{det} A)^{-1}$, i.e., $(\operatorname{det} A)^{2}=1$. This shows that $\operatorname{det} A= \pm 1$, for all $A \in \operatorname{Sp}(n ; \mathbb{R})$. In fact, $\operatorname{det} A=1$ for all $A \in \operatorname{Sp}(n ; \mathbb{R})$, although this is not obvious.
\end{EXA}

\begin{EXA}{LIE-H15-01-28}{$\Sp(2n,\C)$ ist MLG}
One can define a bilinear form $\omega$ on $\mathbb{C}^{2 n}$ by the same formula as in (1.7) (with no conjugates). Over $\mathbb{C}$, we have the relation
$$
\omega(z, w)=(z, \Omega w),
$$
where ( $\cdot, \cdot$ ) is the complex bilinear form in (1.4). The set of $2 n \times 2 n$ complex matrices which preserve this form is the complex symplectic group $\operatorname{Sp}(n ; \mathbb{C})$. A $2 n \times 2 n$ complex matrix $A$ is in $\operatorname{Sp}(n ; \mathbb{C})$ if and only if
$$
-\Omega A^{t r} \Omega=A^{-1} .
$$
holds. (Note: This condition involves $A^{t r}$, not $A^{*}$.) Again, we can easily show that each $A \in \operatorname{Sp}(n ; \mathbb{C})$ satisfies $\operatorname{det} A= \pm 1$ and, again, it is actually the case that $\operatorname{det} A=1$.
\end{EXA}

\begin{EXA}{LIE-H15-01-29}{$E$ als Spezialfall}
The Euclidean group $\mathrm{E}(n)$ is the group of all transformations of $\mathbb{R}^{n}$ that can be expressed as a composition of a translation and an orthogonal linear transformation. We write elements of $\mathrm{E}(n)$ as pairs $\{x, R\}$ with $x \in \mathbb{R}^{n}$ and $R \in \mathrm{O}(n)$, and we let $\{x, R\}$ act on $\mathbb{R}^{n}$ by the formula
$$
\{x, R\} y=R y+x .
$$
Since
$$
\left\{x_{1}, R_{1}\right\}\left\{x_{2}, R_{2}\right\} y=R_{1}\left(R_{2} y+x_{2}\right)+x_{1}=R_{1} R_{2} y+\left(x_{1}+R_{1} x_{2}\right),
$$
the product operation for $\mathrm{E}(n)$ is the following:
$$
\left\{x_{1}, R_{1}\right\}\left\{x_{2}, R_{2}\right\}=\left\{x_{1}+R_{1} x_{2}, R_{1} R_{2}\right\} .
$$
The inverse of an element of $\mathrm{E}(n)$ is given by
$$
\{x, R\}^{-1}=\left\{-R^{-1} x, R^{-1}\right\}
$$
The group $\mathrm{E}(n)$ is not a subgroup of $\mathrm{GL}(n ; \mathbb{R})$, since translations are not linear maps.

However, $\mathrm{E}(n)$ is isomorphic to the (closed) subgroup of $\mathrm{GL}(n+1 ; \mathbb{R})$ consisting of matrices of the form
$$
\left(\begin{array}{ccc} 
& & x_{1} \\
& R & \vdots \\
& & x_{n} \\
0 & \cdots & 0
\end{array}\right)
$$
with $R \in \mathrm{O}(n)$. (The reader may easily verify that matrices of the form (1.11) multiply according to the formula in (1.10).)
\end{EXA}

\begin{EXA}{LIE-H15-01-30}{$P(n;1)$ als Spezialfall}
We similarly define the Poincaré group $\mathrm{P}(n ; 1)$ (also known as the inhomogeneous Lorentz group) to be the group of all transformations of $\mathbb{R}^{n+1}$ of the form
$$
T=T_{x} A
$$
with $x \in \mathbb{R}^{n+1}$ and $A \in \mathrm{O}(n ; 1)$. This group is isomorphic to the group of $(n+2) \times$ ( $n+2$ ) matrices of the form
$$
\left(\begin{array}{ccc} 
& & x_{1} \\
& A & \vdots \\
& & x_{n+1} \\
0 & \cdots & 0
\end{array}\right)
$$
with $A \in \mathrm{O}(n ; 1)$.
\end{EXA}

\begin{EXA}{LIE-H15-01-31}{$H$ als Spezialfall}
The set of all $3 \times 3$ real matrices $A$ of the form
$$
A=\left(\begin{array}{lll}
1 & a & b \\
0 & 1 & c \\
0 & 0 & 1
\end{array}\right)
$$
where $a, b$, and $c$ are arbitrary real numbers, is the Heisenberg group. It is easy to check that the product of two matrices of the form (1.13) is again of that form, and, clearly, the identity matrix is of the form (1.13). Furthermore, direct computation shows that if $A$ is as in (1.13), then
$$
A^{-1}=\left(\begin{array}{ccc}
1 & -a & a c-b \\
0 & 1 & -c \\
0 & 0 & 1
\end{array}\right)
$$
Thus, $H$ is a subgroup of $\operatorname{GL}(3 ; \mathbb{R})$. The Heisenberg group is a model for the Heisenberg-Weyl commutation relations in physics and also serves as a illuminating example for the Baker-Campbell-Hausdorff formula.
\end{EXA}

\begin{EXA}{LIE-H15-01-32}{$H$ als Spezialfall}
The set of all $3 \times 3$ real matrices $A$ of the form
$$
A=\left(\begin{array}{lll}
1 & a & b \\
0 & 1 & c \\
0 & 0 & 1
\end{array}\right)
$$
where $a, b$, and $c$ are arbitrary real numbers, is the Heisenberg group. It is easy to check that the product of two matrices of the form (1.13) is again of that form, and, clearly, the identity matrix is of the form (1.13). Furthermore, direct computation shows that if $A$ is as in (1.13), then
$$
A^{-1}=\left(\begin{array}{ccc}
1 & -a & a c-b \\
0 & 1 & -c \\
0 & 0 & 1
\end{array}\right)
$$
Thus, $H$ is a subgroup of $\operatorname{GL}(3 ; \mathbb{R})$. The Heisenberg group is a model for the Heisenberg-Weyl commutation relations in physics and also serves as a illuminating example for the Baker-Campbell-Hausdorff formula.
\end{EXA}

\begin{EXA}{LIE-H15-01-33}{Zu MLG isomorphe Gruppen}
Several important groups which are not defined as groups of matrices can be thought of as such. The group $\mathbb{R}^{*}$ of non-zero real numbers under multiplication is isomorphic to $\mathrm{GL}(1 ; \mathbb{R})$. Similarly, the group $\mathbb{C}^{*}$ of nonzero complex numbers under multiplication is isomorphic to $\mathrm{GL}(1 ; \mathbb{C})$ and the group $S^{1}$ of complex numbers with absolute value one is isomorphic to $\mathrm{U}(1)$.

The group $\mathbb{R}$ under addition is isomorphic to $\mathrm{GL}(1 ; \mathbb{R})^{+}$( $1 \times 1$ real matrices with positive determinant) via the map $x \rightarrow\left[e^{x}\right]$. The group $\mathbb{R}^{n}$ (with vector addition) is isomorphic to the group of diagonal real matrices with positive diagonal entries, via the map
$$
\left(x_{1}, \ldots, x_{n}\right) \rightarrow\left(\begin{array}{ccc}
e^{x_{1}} & & 0 \\
& \ddots & \\
0 & & e^{x_{n}}
\end{array}\right)
$$
\end{EXA}

\begin{EXA}{LIE-H15-01-47}{$\Ogroup, \SO(n),U(n),\SU(n)$ als kompakte MLG}
The following groups are compact: $\mathrm{O}(n)$ and $\mathrm{SO}(n), \mathrm{U}(n)$ and $\mathrm{SU}(n)$, and $\mathrm{Sp}(n)$. Each of these groups is easily seen to be closed in $M_{n}(\mathbb{C})$ and each satisfies the bound $\left|A_{j k}\right| \leq 1$, since in each case, the columns of $A \in G$ are required to be unit vectors.
\end{EXA}

\begin{EXA}{LIE-H15-01-48}{Spezialfall $\SL(n)$ für $n=1$ als kompakte MLG}
The special linear group $\operatorname{SL}(n ; \mathbb{R})$, for example, is unbounded (except in the trivial case $n=1$ ), because for all $m$, the matrix
$$
A_{m}=\left(\begin{array}{ccccc}
m & & & & \\
& \frac{1}{m} & & & \\
& & 1 & & \\
& & & \ddots & \\
& & & & 1
\end{array}\right)
$$
has determinant one.
\end{EXA}

\begin{EXA}{LIE-H15-01-78}{Determinante als Lie Gruppen Homomorphismus}
The simplest interesting example of a Lie group homomorphism is the determinant, which is a homomorphism of $\operatorname{GL}(n ; \mathbb{C})$ into $\mathbb{C}^{*}$.
\end{EXA}

\begin{EXA}{LIE-H15-01-79}{LG Homomorphismus von $\R$ in $\SO(n)$}
Another simple example is the map $\Phi: \mathbb{R} \rightarrow \mathrm{SO}(2)$ given by
$$
\Phi(\theta)=\left(\begin{array}{rr}
\cos \theta & -\sin \theta \\
\sin \theta & \cos \theta
\end{array}\right) .
$$
This map is clearly continuous, and calculation (using standard trigonometric identities) shows that it is a homomorphism.
\end{EXA}

\begin{EXA}{LIE-H15-01-85}{Beispiel-Anwendung von LG Homomorphismus zwischen $\SU(2)$ und $\SO(3)$}
Suppose, for example, that $U$ is the matrix
$$
U=\left(\begin{array}{cc}
e^{i \theta / 2} & 0 \\
0 & e^{-i \theta / 2}
\end{array}\right) .
$$
Then by direct calculation, we obtain
$$
U\left(\begin{array}{cc}
x_{1} & x_{2}+i x_{3} \\
x_{2}-i x_{3} & -x_{1}
\end{array}\right) U^{-1}=\left(\begin{array}{cc}
x_{1}^{\prime} & x_{2}^{\prime}+i x_{3}^{\prime} \\
x_{2}^{\prime}-i x_{3}^{\prime} & -x_{1}^{\prime}
\end{array}\right),
$$
where $x_{1}^{\prime}=x_{1}$ and
$$
\begin{aligned}
x_{2}^{\prime}+i x_{3}^{\prime} & =e^{i \theta}\left(x_{2}+i x_{3}\right) \\
& =\left(x_{2} \cos \theta-x_{3} \sin \theta\right)+i\left(x_{2} \sin \theta+x_{3} \cos \theta\right) .
\end{aligned}
$$
In this case, then, $\Phi_{U}$ is a rotation by angle $\theta$ in the ( $x_{2}, x_{3}$ )-plane. Note that even though the diagonal entries of $U$ are $e^{ \pm i \theta / 2}$, the map $\Phi_{U}$ is a rotation by angle $\theta$, not $\theta / 2$.
\end{EXA}

\begin{EXA}{LIE-H15-01-89}{$\R^2\times S^1$ Lie Gruppen }
Let
$$
G=\mathbb{R} \times \mathbb{R} \times S^{1}=\left\{(x, y, u) \mid x \in \mathbb{R}, y \in \mathbb{R}, u \in S^{1} \subset \mathbb{C}\right\}
$$
equipped with the group product given by
$$
\left(x_{1}, y_{1}, u_{1}\right) \cdot\left(x_{2}, y_{2}, u_{2}\right)=\left(x_{1}+x_{2}, y_{1}+y_{2}, e^{i x_{1} y_{2}} u_{1} u_{2}\right) .
$$
Then $G$ is a Lie group.

It is easily checked that this operation is associative; the product of three elements with either grouping is
$$
\left(x_{1}+x_{2}+x_{3}, y_{1}+y_{2}+y_{3}, e^{i\left(x_{1} y_{2}+x_{1} y_{3}+x_{2} y_{3}\right)} u_{1} u_{2} u_{3}\right) .
$$
There is an identity element in $G$, namely $e=(0,0,1)$ and each element ( $x, y, u$ ) has an inverse given by $\left(-x,-y, e^{i x y} u^{-1}\right)$. Thus, $G$ is, in fact, a group. Furthermore, both the group product and the map that sends each element to its inverse are clearly smooth, showing that $G$ is a Lie group.
\end{EXA}

\subsection{Exponentiale}

\begin{EXA}{LIE-H15-02-16}{Berechnung von Matrix-Exponentialabbildungen}
Consider the matrices
$$
X_{1}=\left(\begin{array}{rr}
0 & -a \\
a & 0
\end{array}\right) ; \quad X_{2}=\left(\begin{array}{ccc}
0 & a & b \\
0 & 0 & c \\
0 & 0 & 0
\end{array}\right) ; \quad X_{3}=\left(\begin{array}{cc}
a & b \\
0 & a
\end{array}\right)
$$
Then
$$
e^{X_{1}}=\left(\begin{array}{rr}
\cos a & -\sin a \\
\sin a & \cos a
\end{array}\right)
$$
and
$$
e^{X_{2}}=\left(\begin{array}{ccc}
1 & a & b+a c / 2 \\
0 & 1 & c \\
0 & 0 & 1
\end{array}\right)
$$
and
$$
e^{X_{3}}=\left(\begin{array}{cc}
e^{a} & e^{a} b \\
0 & e^{a}
\end{array}\right)
$$

Proof. The eigenvectors of $X_{1}$ are ( $1, i$ ) and ( $i, 1$ ), with eigenvalues $-i a$ and $i a$, respectively. Thus,
$$
e^{X_{1}}=\left(\begin{array}{cc}
1 & i \\
i & 1
\end{array}\right)\left(\begin{array}{cc}
e^{-i a} & 0 \\
0 & e^{i a}
\end{array}\right)\left(\begin{array}{rr}
1 / 2 & -i / 2 \\
-i / 2 & 1 / 2
\end{array}\right)
$$
which simplifies to the claimed result. Meanwhile, $X_{2}^{2}$ has the value $a c$ in the upper right-hand corner and all other entries equal to zero, whereas $X_{2}^{3}=0$. Thus, $e^{X_{2}}= 1+X_{2}+X_{2}^{2} / 2$, which reduces to the claimed result. Finally,
$$
X_{3}=\left(\begin{array}{cc}
a & 0 \\
0 & a
\end{array}\right)+\left(\begin{array}{ll}
0 & b \\
0 & 0
\end{array}\right)
$$
where the two terms clearly commute and the second term is nilpotent. Thus, we obtain
$$
e^{X_{3}}=\left(\begin{array}{cc}
e^{a} & 0 \\
0 & e^{a}
\end{array}\right)\left(\begin{array}{ll}
1 & b \\
0 & 1
\end{array}\right)
$$
which reduces to the claimed result.
\end{EXA}

\subsection{Lie Algebren}

\begin{EXA}{LIE-H15-03-02}{$\R^3$ mit Kreuzprodukt}
Let $\mathfrak{g}=\mathbb{R}^{3}$ and let $[\cdot, \cdot]: \mathbb{R}^{3} \times \mathbb{R}^{3} \rightarrow \mathbb{R}^{3}$ be given by
$$
[x, y]=x \times y,
$$
where $x \times y$ is the cross product (or vector product). Then $\mathfrak{g}$ is a Lie algebra.

Proof. Bilinearity and skew symmetry are standard properties of the cross product. To verify the Jacobi identity, it suffices (by bilinearity) to verify it when $x=e_{j}$, $y=e_{k}$, and $z=e_{l}$, where $e_{1}, e_{2}$, and $e_{3}$ are the standard basis elements for $\mathbb{R}^{3}$. If $j, k$, and $l$ are all equal, each term in the Jacobi identity is zero. If $j, k$, and $l$ are all different, the cross product of any two of $e_{j}, e_{k}$, and $e_{l}$ is equal to a multiple of the third, so again, each term in the Jacobi identity is zero. It remains to consider the case in which two of $j, k, l$ are equal and the third is different. By re-ordering the terms in the Jacobi identity as necessary, it suffices to verify the identity
$$
\left[e_{j},\left[e_{j}, e_{k}\right]\right]+\left[e_{j},\left[e_{k}, e_{j}\right]\right]+\left[e_{k},\left[e_{j}, e_{j}\right]\right] .
$$
The first two terms in (3.1) are negatives of each other and the third is zero.
\end{EXA}

\begin{EXA}{LIE-H15-03-03}{Lie Algebra als Unteralgebra von assoziativen Algebren}
Let $\mathcal{A}$ be an associative algebra and let $\mathfrak{g}$ be a subspace of $\mathcal{A}$ such that $X Y-Y X \in \mathfrak{g}$ for all $X, Y \in \mathfrak{g}$. Then $\mathfrak{g}$ is a Lie algebra with bracket operation given by
$$
[X, Y]=X Y-Y X
$$

Proof. The bilinearity and skew symmetry of the bracket are evident. To verify the Jacobi identity, note that each double bracket generates four terms, for a total of 12 terms. It is left to the reader to verify that the product of $X, Y$, and $Z$ in each of the six possible orderings occurs twice, once with a plus sign and once with a minus sign.
\end{EXA}

\begin{EXA}{LIE-H15-03-05}{$\operatorname{sl}(n;\C)$ (Spurfreie) als Lie Algebra}
Let $\operatorname{sI}(n ; \mathbb{C})$ the space $X \in M_{n}(\mathbb{C})$ for which trace $(X)=0$. Then $\operatorname{sl}(n ; \mathbb{C})$ is a Lie algebra with bracket $[X, Y]=X Y-Y X$.

Proof. For any $X$ and $Y$ in $M_{n}(\mathbb{C})$, we have

$$
\operatorname{trace}(X Y-Y X)=\operatorname{trace}(X Y)-\operatorname{trace}(Y X)=0 .
$$

This holds, in particular, if $X$ and $Y$ have trace zero. Thus, Example 3.3 applies.
\end{EXA}

\begin{EXA}{LIE-H15-03-21}{1-dim Lie Algebra sind nicht einfach und irreduzibel}
A one-dimensional Lie algebra is certainly irreducible, since it is has no nontrivial subspaces and therefore no nontrivial subalgebras and no nontrivial ideals. Nevertheless, such a Lie algebra is, by definition, not considered simple.
\end{EXA}

\begin{EXA}{LIE-H15-03-23}{$\operatorname{sl}(2;\C)$ ist einfach}
The Lie algebra $\mathrm{sl}(2 ; \mathbb{C})$ is simple.

Proof. 

We use the following basis for $\operatorname{sI}(2 ; \mathbb{C})$ :
$$
X=\left(\begin{array}{ll}
0 & 1 \\
0 & 0
\end{array}\right) ; \quad Y=\left(\begin{array}{ll}
0 & 0 \\
1 & 0
\end{array}\right) ; \quad H=\left(\begin{array}{rr}
1 & 0 \\
0 & -1
\end{array}\right)
$$
Direct calculation shows that these basis elements have the following commutation relations: $[X, Y]=H,[H, X]=2 X$, and $[H, Y]=-2 Y$. Suppose $\mathfrak{h}$ is an ideal in $\mathrm{sl}(2 ; \mathbb{C})$ and that $\mathfrak{h}$ contains an element $Z=a X+b H+c Y$, where $a, b$, and $c$ are not all zero. We will show, then, that $\mathfrak{h}=\mathrm{sl}(2 ; \mathbb{C})$. Suppose first that $c \neq 0$. Then the element
$$
[X,[X, Z]]=[X,[-2 b X+c H]]=-2 c X
$$
is a nonzero multiple of $X$. Since $\mathfrak{h}$ is an ideal, we conclude that $X \in \mathfrak{h}$. But $[Y, X]$ is a nonzero multiple of $H$ and $[Y,[Y, X]]$ is a nonzero multiple of $Y$, showing that $Y$ and $H$ also belong to $\mathfrak{h}$, from which we conclude that $\mathfrak{h}=\mathrm{sl}(2 ; \mathbb{C})$.

Suppose next that $c=0$ but $b \neq 0$. Then $[X, Z]$ is a nonzero multiple of $X$ and we may then apply the same argument in the previous paragraph to show that $\mathfrak{h}=\mathrm{sl}(2 ; \mathbb{C})$. Finally, if $c=0$ and $b=0$ but $a \neq 0$, then $Z$ itself is a nonzero multiple of $X$ and we again conclude that $\mathfrak{h}=\mathrm{sl}(2 ; \mathbb{C})$.
\end{EXA}

\begin{EXA}{LIE-H15-03-31}{Obere reelle 3x3 Dreiecksmatrizen mit 0 Diagonale sind nilpotente Lie Algebren}
If $\mathfrak{g} \subset M_{3}(\mathbb{R})$ denotes the space of $3 \times 3$ upper triangular matrices with zeros on the diagonal, then $\mathfrak{g}$ satisfies the assumptions of Example 3.3. The Lie algebra $\mathfrak{g}$ is a nilpotent Lie algebra.

Proof. We will use the following basis for $\mathfrak{g}$,

$$
X=\left(\begin{array}{ccc}
0 & 1 & 0 \\
0 & 0 & 0 \\
0 & 0 & 0
\end{array}\right) ; \quad Y=\left(\begin{array}{ccc}
0 & 0 & 0 \\
0 & 0 & 1 \\
0 & 0 & 0
\end{array}\right) ; \quad Z=\left(\begin{array}{ccc}
0 & 0 & 1 \\
0 & 0 & 0 \\
0 & 0 & 0
\end{array}\right)
$$
Direct calculation then establishes the following commutation relations: $[X, Y]= Z$ and $[X, Z]=[Y, Z]=0$. In particular, the bracket of two elements of $\mathfrak{g}$ is again in $\mathfrak{g}$, so that $\mathfrak{g}$ is a Lie algebra. Then $[\mathfrak{g}, \mathfrak{g}]$ is the span of $Z$ and $[\mathfrak{g},[\mathfrak{g}, \mathfrak{g}]]=0$, showing that $\mathfrak{g}$ is nilpotent.
\end{EXA}

\begin{EXA}{LIE-H15-03-32}{Obere 2x2 komplexe Dreiecksmatrizen sind lösbar aber nicht nilpotent}
If $\mathfrak{g} \subset M_{2}(\mathbb{C})$ denotes the space of $2 \times 2$ matrices of the form
$$
\left(\begin{array}{ll}
a & b \\
0 & c
\end{array}\right)
$$
with $a, b$, and $c$ in $\mathbb{C}$, then $\mathfrak{g}$ satisfies the assumptions of Example 3.3. The Lie algebra $\mathfrak{g}$ is solvable but not nilpotent.

Proof. Direct calculation shows that

$$
\left[\left(\begin{array}{ll}
a & b \\
0 & c
\end{array}\right),\left(\begin{array}{ll}
d & e \\
0 & f
\end{array}\right)\right]=\left(\begin{array}{ll}
0 & h \\
0 & 0
\end{array}\right)
$$
where $h=a e+b f-b d-c e$, showing that $\mathfrak{g}$ is a Lie subalgebra of $M_{2}(\mathbb{C})$. Furthermore, the commutator ideal $[\mathfrak{g}, \mathfrak{g}]$ is one dimensional and hence commutative. Thus, $\mathfrak{g}_{2}=\{0\}$, showing that $\mathfrak{g}$ is solvable. On the other hand, consider the following elements of $\mathfrak{g}$ :
$$
H=\left(\begin{array}{rr}
1 & 0 \\
0 & -1
\end{array}\right) ; \quad X=\left(\begin{array}{ll}
0 & 1 \\
0 & 0
\end{array}\right)
$$
Using (3.6), we can see that $[H, X]=2 X$, and thus that
$$
[H,[H,[H, \cdots[H, X] \cdots]]]
$$
is a nonzero multiple of $X$, showing that $\mathfrak{g}^{j} \neq\{0\}$ for all $j$.
\end{EXA}

\begin{EXA}{LIE-H15-03-41}{$\operatorname{GL}(n ; \mathbb{C}), \operatorname{SL}(n ; \mathbb{C}), \operatorname{SO}(n ; \mathbb{C})$, and $\operatorname{Sp}(n ; \mathbb{C})$ sind komplexe Matrix Lie Gruppen}
Examples of complex groups are $\operatorname{GL}(n ; \mathbb{C}), \operatorname{SL}(n ; \mathbb{C}), \operatorname{SO}(n ; \mathbb{C})$, and $\operatorname{Sp}(n ; \mathbb{C})$, as the calculations in Sect. 3.4 will show.
\end{EXA}

\begin{EXA}{LIE-H15-03-46}{Lie Algebren $\operatorname{gl}(n ; \mathbb{C}), \operatorname{gl}(n ; \mathbb{R}), \operatorname{sl}(n ; \mathbb{C}), \mathrm{sl}(n ; \mathbb{R})$}
The Lie algebra of $\mathrm{GL}(n ; \mathbb{C})$ is the space $M_{n}(\mathbb{C})$ of all $n \times n$ matrices with complex entries. Similarly, the Lie algebra of $\mathrm{GL}(n ; \mathbb{R})$ is equal to $M_{n}(\mathbb{R})$. The Lie algebra of $\operatorname{SL}(n ; \mathbb{C})$ consists of all $n \times n$ complex matrices with trace zero, and the Lie algebra of $\operatorname{SL}(n ; \mathbb{R})$ consists of all $n \times n$ real matrices with trace zero.

We denote the Lie algebras of these groups as $\operatorname{gl}(n ; \mathbb{C}), \operatorname{gl}(n ; \mathbb{R}), \operatorname{sl}(n ; \mathbb{C})$, and $\mathrm{sl}(n ; \mathbb{R})$, respectively.

If $X \in M_{n}(\mathbb{C})$, then $e^{t X}$ is invertible, so that $X$ belongs to the Lie algebra of $\operatorname{GL}(n ; \mathbb{C})$. If $X \in M_{n}(\mathbb{R})$, then $e^{t X}$ is invertible and real, so that $X$ is in the Lie algebra of $\mathrm{GL}(n ; \mathbb{R})$. Conversely, if $e^{t X}$ is real for all real $t$, then $X= d\left(e^{t X}\right) /\left.d t\right|_{t=0}$ must also real. If $X \in M_{n}(\mathbb{C})$ has trace zero, then by Theorem 2.12, $\operatorname{det}\left(e^{t X}\right)=1$, showing that $X$ is in the Lie algebra of $\operatorname{SL}(n ; \mathbb{C})$. Conversely, if $\operatorname{det}\left(e^{t X}\right)=e^{t \operatorname{trace}(X)}=1$ for all real $t$, then
$$
\operatorname{trace}(X)=\left.\frac{d}{d t} e^{t \operatorname{trace}(X)}\right|_{t=0}=0
$$
Finally, if $X$ is real and has trace zero, then $e^{t X}$ is real and has determinant 1 for all real $t$, showing that $X$ is in the Lie algebra of $\operatorname{SL}(n ; \mathbb{R})$. Conversely, if $e^{t X}$ is real and has determinant 1 for all real $t$, the preceding arguments show that $X$ must be real and have trace zero.
\end{EXA}

\begin{EXA}{LIE-H15-03-48}{Lie Algebra $\mathrm{u}(n),\mathrm{su}(n),\mathrm{u}(n),\mathrm{so}(n)$}
The Lie algebra of $U(n)$ consists of all complex matrices satisfying $X^{*}=-X$ and the Lie algebra of $\mathrm{SU}(n)$ consists of all complex matrices satisfying $X^{*}=-X$ and trace $(X)=0$. The Lie algebra of the orthogonal group $\mathrm{O}(n)$ consists of all real matrices $X$ satisfying $X^{t r}=-X$ and the Lie algebra of $\mathrm{SO}(n)$ is the same as that of $\mathrm{O}(n)$.

The Lie algebras of $\mathrm{U}(n)$ and $\mathrm{SU}(n)$ are denoted $\mathrm{u}(n)$ and $\mathrm{su}(n)$, respectively. The Lie algebra of $\mathrm{SO}(n)$ (which is the same as that of $\mathrm{O}(n)$ ) is denoted $\operatorname{so}(n)$.

A matrix $U$ is unitary if and only if $U^{*}=U^{-1}$. Thus, $e^{t X}$ is unitary if and only if
$$
\left(e^{t X}\right)^{*}=\left(e^{t X}\right)^{-1}=e^{-t X}
$$
By Point 2 of Proposition 2.3, $\left(e^{t X}\right)^{*}=e^{t X^{*}}$, and so (3.8) becomes
$$
e^{t X^{*}}=e^{-t X}
$$
The condition (3.9) holds for all real $t$ if and only if $X^{*}=-X$. Thus, the Lie algebra of $\mathrm{U}(n)$ consists precisely of matrices $X$ such that $X^{*}=-X$. As in the proof of Proposition 3.23, adding the "determinant 1" condition at the group level adds the "trace 0 " condition at the Lie algebra level.

An exactly similar argument over $\mathbb{R}$ shows that a real matrix $X$ belongs to the Lie algebra of $\mathrm{O}(n)$ if and only if $X^{t r}=-X$. Since any such matrix has trace $(X)=0$ (since the diagonal entries of $X$ are all zero), we see that every element of the Lie algebra of $\mathrm{O}(n)$ is also in the Lie algebra of $\mathrm{SO}(n)$.
\end{EXA}

\begin{EXA}{LIE-H15-03-53}{Lie Algebra der Heisenberg Gruppe $H$}
The Lie algebra of the Heisenberg group $H$ is the space of all matrices of the form
$$
X=\left(\begin{array}{lll}
0 & a & b \\
0 & 0 & c \\
0 & 0 & 0
\end{array}\right),
$$
with $a, b, c \in \mathbb{R}$.

If $X$ is strictly upper triangular, it is easy to verify that $X^{m}$ will be strictly upper triangular for all positive integers $m$. Thus, for $X$ as in (3.10), we will have $e^{t X}=I+B$ with $B$ strictly upper triangular, showing that $e^{t X} \in H$. Conversely, if $e^{t X}$ belongs to $H$ for all real $t$, then all of the entries of $e^{t X}$ on or below the diagonal are independent of $t$. Thus, $X=d\left(e^{t X}\right) /\left.d t\right|_{t=0}$ will be of the form in (3.10).
\end{EXA}

\begin{EXA}{LIE-H15-03-55}{Basis von $\mathrm{su}(2)$}
The following elements form a basis for the Lie algebra su(2):
$$
E_{1}=\frac{1}{2}\left(\begin{array}{rr}
i & 0 \\
0 & -i
\end{array}\right) ; \quad E_{2}=\frac{1}{2}\left(\begin{array}{ll}
0 & i \\
i & 0
\end{array}\right) ; \quad E_{3}=\frac{1}{2}\left(\begin{array}{rr}
0 & -1 \\
1 & 0
\end{array}\right) .
$$
These elements satisfy the commutation relations $\left[E_{1}, E_{2}\right]=E_{3},\left[E_{2}, E_{3}\right]=E_{1}$, and $\left[E_{3}, E_{1}\right]=E_{2}$.
\end{EXA}

\begin{EXA}{LIE-H15-03-56}{Basis von $\mathrm{so}(3)$}
The following elements form a basis for the Lie algebra so(3):
$$
F_{1}=\left(\begin{array}{rrr}
0 & 0 & 0 \\
0 & 0 & -1 \\
0 & 1 & 0
\end{array}\right), \quad F_{2}=\left(\begin{array}{rrr}
0 & 0 & 1 \\
0 & 0 & 0 \\
-1 & 0 & 0
\end{array}\right), \quad F_{3}=\left(\begin{array}{rrr}
0 & -1 & 0 \\
1 & 0 & 0 \\
0 & 0 & 0
\end{array}\right),
$$
These elements satisfy the commutation relations $\left[F_{1}, F_{2}\right]=F_{3},\left[F_{2}, F_{3}\right]=F_{1}$, and $\left[F_{3}, F_{1}\right]=F_{2}$.
\end{EXA}

\begin{EXA}{LIE-H15-03-60}{Induzierter LA Homomorphismus zwischen $\SU(2)$ und $\SO(3)$}
Let $\Phi: \mathrm{SU}(2) \rightarrow \mathrm{SO}(3); U\mapsto \Phi_{U}(X)=U X U^{-1}$ be the homomorphism in Proposition 1.19. Then the associated Lie algebra homomorphism $\phi: \operatorname{su}(2) \rightarrow \operatorname{so}(3)$ satisfies
$$
\phi\left(E_{j}\right)=F_{j}, \quad j=1,2,3,
$$
where $\left\{E_{1}, E_{2}, E_{3}\right\}$ and $\left\{F_{1}, F_{2}, F_{3}\right\}$ are the bases for $\mathrm{su}(2)$ and so(3), respectively, given in Example 3.27.

Since $\phi$ maps a basis for su(2) to a basis for so(3), we see that $\phi$ is a Lie algebra isomorphism, even though $\Phi$ is not a Lie group isomorphism (since $\operatorname{ker}(\Phi)=\{I,-I\})$.

Proof. If $X$ is in su(2) and $Y$ is in the space $V$ in (1.14)
$$
X=\left(\begin{array}{cc}
x_{1} & x_{2}+i x_{3} \\
x_{2}-i x_{3} & -x_{1}
\end{array}\right),
$$
then
$$
\left.\frac{d}{d t} \Phi\left(e^{t X}\right) Y\right|_{t=0}=\left.\frac{d}{d t} e^{t X} Y e^{-t X}\right|_{t=0}=[X, Y]
$$
Thus, $\phi(X)$ is the linear map of $V \cong \mathbb{R}^{3}$ to itself given by $Y \mapsto[X, Y]$. If, say, $X=E_{1}$, then direct computation shows that
$$
\left[E_{1},\left(\begin{array}{cc}
x_{1} & x_{2}+i x_{3} \\
x_{2}-i x_{3} & -x_{1}
\end{array}\right)\right]=\left(\begin{array}{cc}
x_{1}^{\prime} & x_{2}^{\prime}+i x_{3}^{\prime} \\
x_{2}^{\prime}-i x_{3}^{\prime} & -x_{1}^{\prime}
\end{array}\right),
$$
where $\left(x_{1}^{\prime}, x_{2}^{\prime}, x_{3}^{\prime}\right)=\left(0,-x_{3}, x_{2}\right)$. Since
$$
\left(\begin{array}{c}
0 \\
-x_{3} \\
x_{2}
\end{array}\right)=\left(\begin{array}{rrr}
0 & 0 & 0 \\
0 & 0 & -1 \\
0 & 1 & 0
\end{array}\right)\left(\begin{array}{l}
x_{1} \\
x_{2} \\
x_{3}
\end{array}\right),
$$
we conclude that $\phi\left(E_{1}\right)$ is the $3 \times 3$ matrix appearing on the right-hand side of (3.13), which is precisely $F_{1}$. The computation of $\phi\left(E_{2}\right)$ and $\phi\left(E_{3}\right)$ is similar and is left to the reader.
\end{EXA}

\begin{EXA}{LIE-H15-03-79}{Komplexifizierung von kanonischen Lie Algebren}
Using the proposition, we easily obtain the following list of isomorphisms:

$$
\begin{aligned}
\mathrm{gl}(n ; \mathbb{R})_{\mathbb{C}} & \cong \mathrm{gl}(n ; \mathbb{C}), \\
\mathrm{u}(n)_{\mathbb{C}} & \cong \mathrm{gl}(n ; \mathbb{C}), \\
\mathrm{su}(n)_{\mathbb{C}} & \cong \mathrm{sl}(n ; \mathbb{C}), \\
\mathrm{sl}(n ; \mathbb{R})_{\mathbb{C}} & \cong \mathrm{sl}(n ; \mathbb{C}), \\
\mathrm{so}(n)_{\mathbb{C}} & \cong \mathrm{so}(n ; \mathbb{C}), \\
\mathrm{sp}(n ; \mathbb{R})_{\mathbb{C}} & \cong \mathrm{sp}(n ; \mathbb{C}), \\
\mathrm{sp}(n)_{\mathbb{C}} & \cong \mathrm{sp}(n ; \mathbb{C}) .
\end{aligned}
$$
\end{EXA}

\begin{EXA}{LIE-H15-03-80}{Komplexifizierung von $\operatorname{u}(n)$}
Let us verify just one example, that of $\mathrm{u}(n)$. If $X^{*}=-X$, then $(i X)^{*}=i X$. Thus, $X$ and $X^{*}$ cannot both be in $\mathrm{u}(n)$ unless $X$ is zero. Furthermore, every $X$ in $M_{n}(\mathbb{C})$ can be expressed as $X=X_{1}+i X_{2}$, where $X_{1}=\left(X-X^{*}\right) / 2$ and $X_{2}=\left(X+X^{*}\right) /(2 i)$ are both in $\mathrm{u}(n)$. This shows that $\mathrm{u}(n)_{\mathbb{C}} \cong \mathrm{gl}(n ; \mathbb{C})$.
\end{EXA}

\begin{EXA}{LIE-H15-03-84}{Nicht jedes Exponential ist auch in der Lie Algebra}
There does not exist a matrix $X \in \mathrm{sI}(2 ; \mathbb{C})$ with
$$
e^{X}=\left(\begin{array}{rr}
-1 & 1 \\
0 & -1
\end{array}\right),
$$
even though the matrix on the right-hand side of (3.18) is in $\operatorname{SL}(2 ; \mathbb{C})$.


Proof. If $X \in \operatorname{sl}(2 ; \mathbb{C})$ has distinct eigenvalues, then $X$ is diagonalizable and $e^{X}$ will also be diagonalizable, unlike the matrix on the right-hand side of (3.18). If $X \in \operatorname{sl}(2 ; \mathbb{C})$ has a repeated eigenvalue, this eigenvalue must be 0 or the trace of $X$ would not be zero. Thus, there is a nonzero vector $v$ with $X v=0$, from which it follows that $e^{X} v=e^{0} v=v$. We conclude that $e^{X}$ has 1 as an eigenvalue, unlike the matrix on the right-hand side of (3.18).
\end{EXA}

\section{Propositionen}

\subsection{Matrix Lie Gruppen}

\begin{PROOF}{LIE-H15-01-12}{$\SL(2 ; \mathbb{R})$ ist Gruppe und glatte Untermannigfaltigkeit von $\R^4$}
Since the determinant of a product is the product of the determinants, this set forms a group under the operation of matrix multiplication. If we think of the set of all $2 \times 2$ matrices, with entries $a, b, c, d$, as $\mathbb{R}^{4}$, then $\operatorname{SL}(2 ; \mathbb{R})$ is the set of points in $\mathbb{R}^{4}$ for which the smooth function $a d-b c$ has the value 1 .

Suppose $f$ is a smooth function on $\mathbb{R}^{k}$ and we consider the set $E$ where $f(\mathbf{x})$ equals some constant value $c$. If, at each point $\mathbf{x}_{0}$ in $E$, at least one of the partial derivatives of $f$ is nonzero, then the implicit function theorem tells us that we can solve the equation $f(\mathbf{x})=c$ near $\mathbf{x}_{0}$ for one of the variables as a function of the other $k-1$ variables. Thus, $E$ is a smooth "surface" (or embedded submanifold) in $\mathbb{R}^{k}$ of dimension $k-1$. In the case of $\operatorname{SL}(2 ; \mathbb{R})$ inside $\mathbb{R}^{4}$, we note that the partial derivatives of $a d-b c$ with respect to $a, b, c$, and $d$ are $d,-c,-b$, and $a$, respectively. Thus, at each point where $a d-b c=1$, at least one of these partial derivatives is nonzero, and we conclude that $\mathrm{SL}(2 ; \mathbb{R})$ is a smooth surface of dimension 3 . Thus, $\mathrm{SL}(2 ; \mathbb{R})$ is a Lie group of dimension 3.
\end{PROOF}

\begin{PROOF}{LIE-H15-01-25}{Charakterisierung von $\Ogroup(n,k)$ durch Kommutatoren}
If $A$ is an $(n+k) \times(n+k)$ real matrix, let $A^{(j)}$ denote the $j$ th column vector of $A$, that is,
$$
A^{(j)}=\left(\begin{array}{c}
A_{1, j} \\
\vdots \\
A_{n+k, j}
\end{array}\right) .
$$
Note that $A^{(j)}$ is equal to $A e_{j}$, that is, the result of applying $A$ to the $j$ th standard basis element $e_{j}$. Then $A$ will belong to $\mathrm{O}(n ; k)$ if and only $\left[A e_{j}, A e_{l}\right]$ for all $1 \leq j, l \leq n+k$. Explicitly, this means that $A \in \mathrm{O}(n ; k)$ if and only if the following conditions are satisfied:
$$
\begin{array}{ll}
{\left[A^{(j)}, A^{(l)}\right]_{n, k}=0} & j \neq l, \\
{\left[A^{(j)}, A^{(j)}\right]_{n, k}=1} & 1 \leq j \leq n, \\
{\left[A^{(j)}, A^{(j)}\right]_{n, k}=-1} & n+1 \leq j \leq n+k .
\end{array}
$$
\end{PROOF}

\begin{PROOF}{LIE-H15-01-26}{Charakterisierung von $\Ogroup(n,k)$ durch Einheitsmatrix mit $(n,k)$-Signatur}
Let $g$ denote the $(n+k) \times(n+k)$ diagonal matrix with ones in the first $n$ diagonal entries and minus ones in the last $k$ diagonal entries:
$$
g=\left(\begin{array}{llllll}
1 & & & & & \\
& \ddots & & & & \\
& & 1 & & & \\
& & & -1 & & \\
& & & & \ddots & \\
& & & & & -1
\end{array}\right)
$$
Then $A$ is in $\mathrm{O}(n ; k)$ if and only if $A^{t r} g A=g$. 

Taking the determinant of this equation gives $(\operatorname{det} A)^{2} \operatorname{det} g=\operatorname{det} g$, or $(\operatorname{det} A)^{2}=1$. Thus, for any $A$ in $\mathrm{O}(n ; k), \operatorname{det} A= \pm 1$.
\end{PROOF}

\begin{PROP}{LIE-H15-01-35}{Charakterisierung von $\U(2n)$-Elementen in $\Sp(n)$}
If $\U$ belongs to $\U(2n)$ then $\U$ belongs to $\Sp(n)$ if and only if $\U$ commutes with $J$.
\end{PROP}

\begin{PROOF}{LIE-H15-01-37}{$\Sp(n)$ als Untiere Gruppe über den Quaternionen}
The quaternion algebra $\mathbb{H}$ is the four-dimensional associative algebra over $\mathbb{R}$ spanned by elements 1 (the identity), $\mathbf{i}, \mathbf{j}$, and $\mathbf{k}$ satisfying
$$
\mathbf{i}^{2}=\mathbf{j}^{2}=\mathbf{k}^{2}=-1
$$
and
$$
\begin{aligned}
\mathbf{i j}=\mathbf{k} ; & \mathbf{j i}=-\mathbf{k} ; \\
\mathbf{j k}=\mathbf{i} ; & \mathbf{k j}=-\mathbf{i} ; \\
\mathbf{k i}=\mathbf{j} ; & \mathbf{i k}=-\mathbf{j} .
\end{aligned}
$$
We may realize the quaternion algebra inside $M_{2}(\mathbb{C})$ by taking identifying 1 with the identity matrix and setting
$$
\mathbf{i}=\left(\begin{array}{rr}
i & 0 \\
0 & -i
\end{array}\right) ; \quad \mathbf{j}=\left(\begin{array}{rr}
0 & 1 \\
-1 & 0
\end{array}\right) ; \quad \mathbf{k}=\left(\begin{array}{cc}
0 & i \\
i & 0
\end{array}\right) .
$$
The algebra $\mathbb{H}$ is then the space of real linear combinations of $I, \mathbf{i}, \mathbf{j}$, and $\mathbf{k}$.


Now, since $J$ is conjugate linear, we have
$$
J(i z)=-i J z
$$
for all $z \in \mathbb{C}^{2 n}$; that is, $i J=-J i$. Thus, if we define $K$ to be $i J$, we have
$$
K^{2}=i J i J=-J(i)^{2} J=J^{2}=-I
$$
and one can easily check that $i I, J$, and $K$ satisfy the same commutation relations as $\mathbf{i}$, $\mathbf{j}$, and $\mathbf{k}$. We can therefore make $\mathbb{C}^{2 n}$ into a "vector space" over the noncommutative algebra $\mathbb{H}$ by setting
$$
\begin{aligned}
\mathbf{i} \cdot z & =i z \\
\mathbf{j} \cdot z & =J z \\
\mathbf{k} \cdot z & =i J z .
\end{aligned}
$$
Now, if $U$ belongs to $\operatorname{Sp}(n)$, then $U$ commutes with multiplication by $i$ and with $J$ (Proposition 1.5) and thus, also, with $K:=i J$. Thus, $U$ is actually "quaternion linear." A $2 n \times 2 n$ matrix $U$ therefore belongs to $\mathrm{Sp}(n)$ if and only if $U$ is quaternion linear and preserves the norm. 

Thus, we may think of $\operatorname{Sp}(n)$ as the "unitary group over the quaternions." The compact symplectic group then fits naturally with the orthogonal groups (norm-preserving maps over $\mathbb{R}$ ) and the unitary groups (normpreserving maps over $\mathbb{C}$ ).
\end{PROOF}

\begin{THEO}{LIE-H15-01-38}{Charakterisierung von Eigenwerten/vektoren von $\U(2n)$ in $\Sp(n)$}
If $U \in \operatorname{Sp}(n)$, then there exists an orthonormal basis $u_{1}, \ldots, u_{n}$, $v_{1}, \ldots v_{n}$ for $\mathbb{C}^{n}$ such that the following properties hold: 

\begin{enumerate}
\item $J u_{j}=v_{j}$
\item For some real numbers $\theta_{1}, \ldots, \theta_{n}$, we have
$$
\begin{aligned}
U u_{j} & =e^{i \theta_{j}} u_{j} \\
U v_{j} & =e^{-i \theta_{j}} v_{j} ;
\end{aligned}
$$
\item $$
\begin{aligned}
& \omega\left(u_{j}, u_{k}\right)=\omega\left(v_{j}, v_{k}\right)=0 \\
& \omega\left(u_{j}, v_{k}\right)=\delta_{j k}
\end{aligned}
$$
\end{enumerate}
Conversely, if there exists an orthonormal basis with these properties, $U$ belongs to $\mathrm{Sp}(n)$.
\end{THEO}

\begin{LEM}{LIE-H15-01-40}{Invarianz-Eigenschaften von $J$ auf komplexen Unterräumen}
Suppose $V$ is a complex subspace of $\mathbb{C}^{2 n}$ that is invariant under the conjugate-linear map $J$. Then the orthogonal complement $V^{\perp}$ of $V$ (with respect to the inner product $\langle\cdot, \cdot\rangle$ ) is also invariant under $J$. Furthermore, $V$ and $V^{\perp}$ are orthogonal with respect to $\omega$; that is,
$$
\omega(z, w)=0
$$
for all $z \in V$ and $w \in V^{\perp}$.
\end{LEM}

\begin{LEM}{LIE-H15-01-41}{Invarianz-Eigenschaften von $J$ auf komplexen Unterräumen}
Suppose $V$ is a complex subspace of $\mathbb{C}^{2 n}$ that is invariant under the conjugate-linear map $J$. Then the orthogonal complement $V^{\perp}$ of $V$ (with respect to the inner product $\langle\cdot, \cdot\rangle$ ) is also invariant under $J$. Furthermore, $V$ and $V^{\perp}$ are orthogonal with respect to $\omega$; that is,
$$
\omega(z, w)=0
$$
for all $z \in V$ and $w \in V^{\perp}$.
\end{LEM}

\begin{LEM}{LIE-H15-01-42}{Invarianz-Eigenschaften von $J$ auf komplexen Unterräumen}
Suppose $V$ is a complex subspace of $\mathbb{C}^{2 n}$ that is invariant under the conjugate-linear map $J$. Then the orthogonal complement $V^{\perp}$ of $V$ (with respect to the inner product $\langle\cdot, \cdot\rangle$ ) is also invariant under $J$. Furthermore, $V$ and $V^{\perp}$ are orthogonal with respect to $\omega$; that is,
$$
\omega(z, w)=0
$$
for all $z \in V$ and $w \in V^{\perp}$.
\end{LEM}

\begin{LEM}{LIE-H15-01-43}{Invarianz-Eigenschaften von $J$ auf komplexen Unterräumen}
Suppose $V$ is a complex subspace of $\mathbb{C}^{2 n}$ that is invariant under the conjugate-linear map $J$. Then the orthogonal complement $V^{\perp}$ of $V$ (with respect to the inner product $\langle\cdot, \cdot\rangle$ ) is also invariant under $J$. Furthermore, $V$ and $V^{\perp}$ are orthogonal with respect to $\omega$; that is,
$$
\omega(z, w)=0
$$
for all $z \in V$ and $w \in V^{\perp}$.
\end{LEM}

\begin{PROP}{LIE-H15-01-52}{$id$-WZSH-Komponente ist normale Untergruppe in MLG}
If $G$ is a matrix Lie group, the identity component $G_{0}$ of $G$ is a normal subgroup of $G$.
\end{PROP}

\begin{PROOF}{LIE-H15-01-54}{Stetigkeit von Komposition und Inversion von Wegen in $\GL(n;\C)$}
Note that because matrix multiplication and matrix inversion are continuous on $\mathrm{GL}(n ; \mathbb{C})$, it follows that if $A(t)$ and $B(t)$ are continuous, then so are $A(t) B(t)$ and $A(t)^{-1}$. The continuity of the matrix product is obvious. The continuity of the inverse follows from the formula for the inverse in terms of cofactors; this formula is continuous as long as we remain in the set of invertible matrices where the determinant in the denominator is nonzero.
\end{PROOF}

\begin{PROP}{LIE-H15-01-55}{$\GL(n;\C)$ ist WZSH für $n\geq 1$}
The group $\mathrm{GL}(n ; \mathbb{C})$ is connected for all $n \geq 1$.
\end{PROP}

\begin{PROP}{LIE-H15-01-57}{$\SL(n;\C)$ ist WZSH für $n\geq 1$}
The group $\operatorname{SL}(n ; \mathbb{C})$ is connected for all $n \geq 1$.
\end{PROP}

\begin{PROP}{LIE-H15-01-59}{$\U(n)$ ist WZSH für $n\geq 1$}
The groups $\mathrm{U}(n)$ and $\mathrm{SU}(n)$ are connected, for all $n \geq 1$.
\end{PROP}

\begin{PROP}{LIE-H15-01-61}{$\SU(n)$ ist WZSH für $n\geq 1$}
The groups $\mathrm{SU}(n)$ and $\mathrm{SU}(n)$ are connected, for all $n \geq 1$.
\end{PROP}

\begin{PROP}{LIE-H15-01-63}{$\SO(n)$ ist WZSH für $n\geq 1$}
The group $\mathrm{SO}(n)$ is connected, for all $n \geq 1$.
\end{PROP}

\begin{PROP}{LIE-H15-01-66}{$\SU(2)$ ist EZSH}
The group $\mathrm{SU}(2)$ is simply connected.
\end{PROP}

\begin{PROOF}{LIE-H15-01-69}{$\R P^n$ lokal isomorph zu $S^n$}
There is a natural map $\pi: S^{n} \rightarrow \mathbb{R} P^{n}$, given by
$$
\pi(u)=\{u,-u\} .
$$
We may define a distance function on $\mathbb{R} P^{n}$ by defining
$$
\begin{aligned}
d(\{u,-u\},\{v,-v\}) & =\min (d(u, v), d(u,-v), d(-u, v), d(-u,-v)) \\
& =\min (d(u, v), d(u,-v))
\end{aligned}
$$
(The second equality holds because $d(x, y)=d(-x,-y)$.) With this metric, $\mathbb{R} P^{n}$ is locally isometric to $S^{n}$, since if $u$ and $v$ are nearby points in $S^{n}$, we have $d(\{u,-u\},\{v,-v\})=d(u, v)$.
\end{PROOF}

\begin{PROP}{LIE-H15-01-70}{$\R P^n$ nicht EZSH}
It is known that $\mathbb{R} P^{n}$ is not simply connected. (See, for example, Example 1.43 in [Hat].) Indeed, suppose $u$ is any unit vector in $\mathbb{R}^{n+1}$ and $B(t)$ is any path in $S^{n}$ connecting $u$ to $-u$. Then
$$
A(t):=\pi(B(t))
$$
is a loop in $\mathbb{R} P^{n}$, and this loop cannot be shrunk continuously to a point in $\mathbb{R} P^{n}$. To prove this claim, suppose that a map $A(s, t)$ as in Definition 1.14. Then $A(s, t)$ can be "lifted" to a continuous map $B(s, t)$ into $S^{n}$ such that $B(0, t)=B(t)$ and such that $A(s, t)=\pi(B(s, t))$. (See Proposition 1.30 in [Hat].) Since $A(s, 0)= A(s, 1)$ for all $s$, we must have $B(s, 0)= \pm B(s, 1)$. But by construction, $B(0,0)= -B(0,1)$. If order for $B(s, t)$ to be continuous in $s$, we must then have $B(s, 0)= -B(s, 1)$ for all $s$. It follows that $B(1, t)$ is a nonconstant path in $S^{n}$. It is then easily verified that $A(1, t)=\pi(B(1, t))$ cannot be constant, contradicting our assumption about $A(s, t)$.
\end{PROP}

\begin{PROP}{LIE-H15-01-73}{$\SO(3)$ und $\R P^3$ sind homöomorph und $\SO(3)$ damit nicht EZSH}
There is a continuous bijection between $\mathrm{SO}(3)$ and $\mathbb{R} P^{3}$. Since $\mathbb{R} P^{3}$ is not simply connected, it follows that $\mathrm{SO}(3)$ is not simply connected, either.
\end{PROP}

\begin{PROP}{LIE-H15-01-80}{LG Homomorphismus zwischen $\SU(2)$ und $\SO(3)$}
Die Abbildung definiert durch $\Phi_{U}(X)=U X U^{-1}$ von $\SU(2)$ nach $\SO(3)$ ist ein Lie Gruppen Homomorphismus.
\end{PROP}

\begin{PROP}{LIE-H15-01-81}{LG Homomorphismus zwischen $\SU(2)$ und $\SO(3)$}
Die Abbildung definiert durch $\Phi_{U}(X)=U X U^{-1}$ von $\SU(2)$ nach $\SO(3)$ ist ein Lie Gruppen Homomorphismus.
\end{PROP}

\begin{PROP}{LIE-H15-01-83}{LG Homomorphismus zwischen $\SU(2)$ und $\SO(3)$ ist 2-1 und Onto}
The map $U \mapsto \Phi_{U}$ is a 2-1 and onto map of $\mathrm{SU}(2)$ to $\mathrm{SO}(3)$, with kernel equal to $\{I,-I\}$.
\end{PROP}

\begin{PROOF}{LIE-H15-01-86}{$\SO(3)$ nicht EZSH via Überlagerung}
Since $\mathrm{SU}(2)$ is homeomorphic to $S^{3}$, the proposition gives another way of seeing that SO (3) is homeomorphic to $\mathbb{R} P^{3}$, that is, $S^{3}$ with antipodal points identified. This result was obtained in a different way in Proposition 1.17.

It is not hard to show that $\Phi$ is a "covering map" in the topological sense (Section 1.3 of [Hat]). Since SU(2) is simply connected (Proposition 1.15) and the map is 2-1, it follows by the theory of covering maps (e.g., Theorem 1.38 in [Hat]) that SO(3) cannot be simply connected and, indeed, it must have fundamental group $\mathbb{Z} / 2$. See Chapter 13 for general information about fundamental groups and for a computation of the fundamental group of $\mathrm{SO}(n), n>2$.
\end{PROOF}

\subsection{Exponentiale}

\begin{PROP}{LIE-H15-02-02}{Exponentialabbildung konvergiert überall in $M_n(\C)$}
The series 
$$
e^{X}=\sum_{m=0}^{\infty} \frac{X^{m}}{m!},
$$
converges for all $X \in M_{n}(\mathbb{C})$ and $e^{X}$ is a continuous function of $X$.
\end{PROP}

\begin{PROP}{LIE-H15-02-04}{Ungleichungen der Hilbert-Schmidt-Norm}
The Hilbert-Schmidt norm satisfies the inequalities
$$
\begin{aligned}
\|X+Y\| & \leq\|X\|+\|Y\| \\
\|X Y\| & \leq\|X\|\|Y\|
\end{aligned}
$$
for all $X, Y \in M_{n}(\mathbb{C})$.
\end{PROP}

\begin{PROP}{LIE-H15-02-05}{Ungleichungen der Hilbert-Schmidt-Norm}
The Hilbert-Schmidt norm satisfies the inequalities
$$
\begin{aligned}
\|X+Y\| & \leq\|X\|+\|Y\| \\
\|X Y\| & \leq\|X\|\|Y\|
\end{aligned}
$$
for all $X, Y \in M_{n}(\mathbb{C})$.
\end{PROP}

\begin{PROP}{LIE-H15-02-09}{Elementare Eigenschaften der Exponentialabbildung}
Let $X$ and $Y$ be arbitrary $n \times n$ matrices. Then we have the following:
\begin{enumerate}
  \item $e^{0}=I$.
  \item $\left(e^{X}\right)^{*}=e^{X^{*}}$.
  \item $e^{X}$ is invertible and $\left(e^{X}\right)^{-1}=e^{-X}$.
  \item $e^{(\alpha+\beta) X}=e^{\alpha X} e^{\beta X}$ for all $\alpha$ and $\beta$ in $\mathbb{C}$.
  \item If $X Y=Y X$, then $e^{X+Y}=e^{X} e^{Y}=e^{Y} e^{X}$.
  \item If $C$ is invertible, then $e^{C X C^{-1}}=C e^{X} C^{-1}$.
\end{enumerate}
Although $e^{X+Y}=e^{X} e^{Y}$ when $X$ and $Y$ commute, this identity fails in general.
\end{PROP}

\begin{PROP}{LIE-H15-02-11}{Induzierte glatte Kurve durch Exponentialabbildung}
Let $X$ be a $n \times n$ complex matrix. Then $e^{t X}$ is a smooth curve in $M_{n}(\mathbb{C})$ and
$$
\frac{d}{d t} e^{t X}=X e^{t X}=e^{t X} X
$$
In particular,
$$
\left.\frac{d}{d t} e^{t X}\right|_{t=0}=X
$$
\end{PROP}

\begin{LEM}{LIE-H15-02-18}{Logarithmus in komplexer Analysis}
The function
$$
\log z=\sum_{m=1}^{\infty}(-1)^{m+1} \frac{(z-1)^{m}}{m}
$$
is defined and analytic in a circle of radius 1 about $z=1$.

For all $z$ with $|z-1|<1$,
$$
e^{\log z}=z .
$$
For all $u$ with $|u|<\log 2$, we have $\left|e^{u}-1\right|<1$ and
$$
\log e^{u}=u .
$$
\end{LEM}

\begin{THEO}{LIE-H15-02-22}{Wohldefiniertheit von Log-Abbildung und Bedingungen für Investierung von Exp-Abbildung}
The function
$$
\log A=\sum_{m=1}^{\infty}(-1)^{m+1} \frac{(A-I)^{m}}{m}
$$
is defined and continuous on the set of all $n \times n$ complex matrices $A$ with $\|A-I\|<1$.

For all $A$ with $\|A-I\|<1$,
$$
e^{\log A}=A
$$
For all $X$ with $\|X\|<\log 2,\left\|e^{X}-I\right\|<1$ and
$$
\log e^{X}=X
$$
\end{THEO}

\begin{PROP}{LIE-H15-02-25}{Abschätzung von $\log (I+b)-B$}
There exists a constant $c$ such that for all $n \times n$ matrices $B$ with $\|B\|<\frac{1}{2}$,
$$
\|\log (I+B)-B\| \leq c\|B\|^{2}
$$
\end{PROP}

\begin{THEO}{LIE-H15-02-28}{Darstellung von invertierbaren Matrizen durch Exp-Abbildung}
Every invertible $n \times n$ matrix can be expressed as $e^{X}$ for some $X \in M_{n}(\mathbb{C})$.
\end{THEO}

\begin{THEO}{LIE-H15-02-29}{Lie-Produkt Regel}
For all $X, Y \in M_{n}(\mathbb{C})$, we have
$$
e^{X+Y}=\lim _{m \rightarrow \infty}\left(e^{\frac{X}{m}} e^{\frac{Y}{m}}\right)^{m}
$$
\end{THEO}

\begin{THEO}{LIE-H15-02-31}{Determinante des Exponentialabbildung}
For any $X \in M_{n}(\mathbb{C})$, we have
$$
\operatorname{det}\left(e^{X}\right)=e^{\operatorname{trace}(X)}
$$
\end{THEO}

\begin{THEO}{LIE-H15-02-34}{Darstellungen von 1-Parameter Untergruppen von $\GL(n;\C)$ durch Exp-Abbildung}
If $A(\cdot)$ is a one-parameter subgroup of $\mathrm{GL}(n ; \mathbb{C})$, there exists a unique $n \times n$ complex matrix $X$ such that
$$
A(t)=e^{t X}
$$
\end{THEO}

\begin{LEM}{LIE-H15-02-36}{Eindeutige Existenz von Wurzeln für Elemente im Bild von Exp-Abbildung von hinreichend kleiner Umgebung von Ursprung}
Fix some $\varepsilon$ with $\varepsilon<\log 2$. Let $B_{\varepsilon / 2}$ be the ball of radius $\varepsilon / 2$ around the origin in $M_{n}(\mathbb{C})$, and let $U=\exp \left(B_{\varepsilon / 2}\right)$. Then every $B \in U$ has a unique square root $C$ in $U$, given by $C=\exp \left(\frac{1}{2} \log B\right)$.
\end{LEM}

\begin{PROP}{LIE-H15-02-38}{Exp-Abbildung ist $C^\infty(M_n(\C),M_n(\C))$}
The exponential map is an infinitely differentiable map of $M_{n}(\mathbb{C})$ into $M_{n}(\mathbb{C})$.
\end{PROP}

\begin{THEO}{LIE-H15-02-42}{Eigenschaften von Dekompositionen von invertierbaren Matrizen}
\begin{enumerate}
  \item Every $A \in \mathrm{GL}(n ; \mathbb{C})$ can be written uniquely in the form
    $$
    A=U P
    $$
    where $U$ is unitary and $P$ is self-adjoint and positive.
    
    \item Every self-adjoint positive matrix $P$ can be written uniquely in the form
    $$
    P=e^{X}
    $$
    with $X$ self-adjoint. Conversely, if $X$ is self-adjoint, then $e^{X}$ is self-adjoint and positive.

    \item If we decompose each $A \in \mathrm{GL}(n ; \mathbb{C})$ (uniquely) as
    $$
    A=U e^{X}
    $$
    with $U$ unitary and $X$ self-adjoint, then $U$ and $X$ depend continuously on $A$.
\end{enumerate}
\end{THEO}

\begin{LEM}{LIE-H15-02-44}{Eindeutige selbstadjungierte Wurzel für selbstadjungierte positive Matrizen}
If $Q$ is a self-adjoint, positive matrix, then $Q$ has a unique positive, self-adjoint square root.
\end{LEM}

\begin{PROP}{LIE-H15-02-46}{Polar-Dekomposition von $\GL(n;\R),\SL(n;\C),\SO(n;\R)$}
\begin{enumerate}
\item Every $A \in \mathrm{GL}(n ; \mathbb{R})$ can be written uniquely as
$$
A=\operatorname{Re}^{X}
$$
where $R$ is in $\mathrm{O}(n)$ and $X$ is real and symmetric.\\

\item Every $A \in \mathrm{SL}(n ; \mathbb{C})$ can be written uniquely as
$$
A=U e^{X}
$$
where $U$ is in $\mathrm{SU}(n)$ and $X$ is self-adjoint with trace zero.

\item Every $A \in \mathrm{SL}(n ; \mathbb{R})$ can be written uniquely as
$$
A=R e^{X}
$$
where $R$ is in $\mathrm{SO}(n)$ and $X$ is real and symmetric and has trace zero.
\end{enumerate}
\end{PROP}

\subsection{Lie Algebren}

\begin{PROP}{LIE-H15-03-12}{Adjungierte Abbildung als Lie Algebren Homomorphismus}
If $\mathfrak{g}$ is a Lie algebra, then
$$
\operatorname{ad}_{[X, Y]}=\operatorname{ad}_{X} \operatorname{ad}_{Y}-\operatorname{ad}_{Y} \operatorname{ad}_{X}=\left[\operatorname{ad}_{X}, \operatorname{ad}_{Y}\right] ;
$$
that is, ad: $\mathfrak{g} \rightarrow \operatorname{End}(\mathfrak{g})$ is a Lie algebra homomorphism.
\end{PROP}

\begin{PROP}{LIE-H15-03-35}{$e^X$ von Lie Algebren Elementen sind in Identität-Komponente der Lie Gruppe }
Let $G$ be a matrix Lie group, and $X$ an element of its Lie algebra. Then $e^{X}$ is an element of the identity component $G_{0}$ of $G$.
\end{PROP}

\begin{THEO}{LIE-H15-03-37}{Elementare Eigenschaften von Lie Algebren von Matrix Lie Gruppen}
Let $G$ be a matrix Lie group with Lie algebra $\mathfrak{g}$. If $X$ and $Y$ are elements of $\mathfrak{g}$, the following results hold.
\begin{enumerate}
  \item $A X A^{-1} \in \mathfrak{g}$ for all $A \in G$.
  \item $s X \in \mathfrak{g}$ for all real numbers $s$.
  \item $X+Y \in \mathfrak{g}$.
  \item $X Y-Y X \in \mathfrak{g}$.
\end{enumerate}
It follows from this result and Example 3.3 that the Lie algebra of a matrix Lie group is a real Lie algebra, with bracket given by $[X, Y]=X Y-Y X$. For $X$ and $Y$ in $\mathfrak{g}$, we refer to $[X, Y]=X Y-Y X \in \mathfrak{g}$ as the bracket or commutator of $X$ and $Y$.
\end{THEO}

\begin{PROP}{LIE-H15-03-42}{Lie Algebren von kommutative Matrix Lie Gruppen sind kommutativ}
If $G$ is commutative then $\mathfrak{g}$ is commutative.
\end{PROP}

\begin{PROP}{LIE-H15-03-44}{Lie Algebren $\operatorname{gl}(n ; \mathbb{C}), \operatorname{gl}(n ; \mathbb{R}), \operatorname{sl}(n ; \mathbb{C}), \mathrm{sl}(n ; \mathbb{R})$}
The Lie algebra of $\mathrm{GL}(n ; \mathbb{C})$ is the space $M_{n}(\mathbb{C})$ of all $n \times n$ matrices with complex entries. Similarly, the Lie algebra of $\mathrm{GL}(n ; \mathbb{R})$ is equal to $M_{n}(\mathbb{R})$. The Lie algebra of $\operatorname{SL}(n ; \mathbb{C})$ consists of all $n \times n$ complex matrices with trace zero, and the Lie algebra of $\operatorname{SL}(n ; \mathbb{R})$ consists of all $n \times n$ real matrices with trace zero.

We denote the Lie algebras of these groups as $\operatorname{gl}(n ; \mathbb{C}), \operatorname{gl}(n ; \mathbb{R}), \operatorname{sl}(n ; \mathbb{C})$, and $\mathrm{sl}(n ; \mathbb{R})$, respectively.
\end{PROP}

\begin{PROP}{LIE-H15-03-45}{Lie Algebren $\operatorname{gl}(n ; \mathbb{C}), \operatorname{gl}(n ; \mathbb{R}), \operatorname{sl}(n ; \mathbb{C}), \mathrm{sl}(n ; \mathbb{R})$}
The Lie algebra of $\mathrm{GL}(n ; \mathbb{C})$ is the space $M_{n}(\mathbb{C})$ of all $n \times n$ matrices with complex entries. Similarly, the Lie algebra of $\mathrm{GL}(n ; \mathbb{R})$ is equal to $M_{n}(\mathbb{R})$. The Lie algebra of $\operatorname{SL}(n ; \mathbb{C})$ consists of all $n \times n$ complex matrices with trace zero, and the Lie algebra of $\operatorname{SL}(n ; \mathbb{R})$ consists of all $n \times n$ real matrices with trace zero.

We denote the Lie algebras of these groups as $\operatorname{gl}(n ; \mathbb{C}), \operatorname{gl}(n ; \mathbb{R}), \operatorname{sl}(n ; \mathbb{C})$, and $\mathrm{sl}(n ; \mathbb{R})$, respectively.
\end{PROP}

\begin{PROP}{LIE-H15-03-47}{Lie Algebra $\mathrm{u}(n),\mathrm{su}(n),\mathrm{u}(n),\mathrm{so}(n)$}
The Lie algebra of $U(n)$ consists of all complex matrices satisfying $X^{*}=-X$ and the Lie algebra of $\mathrm{SU}(n)$ consists of all complex matrices satisfying $X^{*}=-X$ and trace $(X)=0$. The Lie algebra of the orthogonal group $\mathrm{O}(n)$ consists of all real matrices $X$ satisfying $X^{t r}=-X$ and the Lie algebra of $\mathrm{SO}(n)$ is the same as that of $\mathrm{O}(n)$.

The Lie algebras of $\mathrm{U}(n)$ and $\mathrm{SU}(n)$ are denoted $\mathrm{u}(n)$ and $\mathrm{su}(n)$, respectively. The Lie algebra of $\mathrm{SO}(n)$ (which is the same as that of $\mathrm{O}(n)$ ) is denoted $\operatorname{so}(n)$.
\end{PROP}

\begin{PROP}{LIE-H15-03-49}{Lie Algebra $\mathrm{so}(n;k),\mathrm{sp}(n;\R),\mathrm{sp}(n,\C),\mathrm{sp}(n)$}
If $g$ is the matrix according to $g A^{t r} g=A^{-1}$, then the Lie algebra of $\mathrm{O}(n ; k)$ consists precisely of those real matrices $X$ such that
$$
g X^{t r} g=-X,
$$
and the Lie algebra of $\mathrm{SO}(n ; k)$ is the same as that of $\mathrm{O}(n ; k)$. If $\Omega$ is the matrix (1.8), then the Lie algebra of $\operatorname{Sp}(n ; \mathbb{R})$ consists precisely of those real matrices $X$ such that
$$
\Omega X^{t r} \Omega=X,
$$
and the Lie algebra of $\operatorname{Sp}(n ; \mathbb{C})$ consists precisely of those complex matrices $X$ satisfying the same condition. The Lie algebra of $\operatorname{Sp}(n)$ consists precisely of those complex matrices $X$ such that $\Omega X^{t r} \Omega=X$ and $X^{*}=-X$.

The Lie algebra of $\mathrm{SO}(n ; k)$ (which is the same as that of $\mathrm{O}(n ; k)$ ) is denoted $\operatorname{so}(n ; k)$, whereas the Lie algebras of the symplectic groups are denoted $\operatorname{sp}(n ; \mathbb{R}), \operatorname{sp}(n ; \mathbb{C})$, and $\operatorname{sp}(n)$.
\end{PROP}

\begin{PROP}{LIE-H15-03-50}{Lie Algebra $\mathrm{so}(n;k),\mathrm{sp}(n;\R),\mathrm{sp}(n,\C),\mathrm{sp}(n)$}
If $g$ is the matrix according to $g A^{t r} g=A^{-1}$, then the Lie algebra of $\mathrm{O}(n ; k)$ consists precisely of those real matrices $X$ such that
$$
g X^{t r} g=-X,
$$
and the Lie algebra of $\mathrm{SO}(n ; k)$ is the same as that of $\mathrm{O}(n ; k)$. If $\Omega$ is the matrix (1.8), then the Lie algebra of $\operatorname{Sp}(n ; \mathbb{R})$ consists precisely of those real matrices $X$ such that
$$
\Omega X^{t r} \Omega=X,
$$
and the Lie algebra of $\operatorname{Sp}(n ; \mathbb{C})$ consists precisely of those complex matrices $X$ satisfying the same condition. The Lie algebra of $\operatorname{Sp}(n)$ consists precisely of those complex matrices $X$ such that $\Omega X^{t r} \Omega=X$ and $X^{*}=-X$.

The Lie algebra of $\mathrm{SO}(n ; k)$ (which is the same as that of $\mathrm{O}(n ; k)$ ) is denoted $\operatorname{so}(n ; k)$, whereas the Lie algebras of the symplectic groups are denoted $\operatorname{sp}(n ; \mathbb{R}), \operatorname{sp}(n ; \mathbb{C})$, and $\operatorname{sp}(n)$.
\end{PROP}

\begin{PROP}{LIE-H15-03-51}{Lie Algebra der Heisenberg Gruppe $H$}
The Lie algebra of the Heisenberg group $H$ is the space of all matrices of the form
$$
X=\left(\begin{array}{lll}
0 & a & b \\
0 & 0 & c \\
0 & 0 & 0
\end{array}\right),
$$
with $a, b, c \in \mathbb{R}$.
\end{PROP}

\begin{PROP}{LIE-H15-03-52}{Lie Algebra der Heisenberg Gruppe $H$}
The Lie algebra of the Heisenberg group $H$ is the space of all matrices of the form
$$
X=\left(\begin{array}{lll}
0 & a & b \\
0 & 0 & c \\
0 & 0 & 0
\end{array}\right),
$$
with $a, b, c \in \mathbb{R}$.
\end{PROP}

\begin{THEO}{LIE-H15-03-57}{Induzierter LA Homomorphismus durch LG Homomorphismus zwischen Matrix Lie Gruppen}
Let $G$ and $H$ be matrix Lie groups, with Lie algebras $\mathfrak{g}$ and $\mathfrak{h}$, respectively. Suppose that $\Phi: G \rightarrow H$ is a Lie group homomorphism. Then there exists a unique real-linear map $\phi: \mathfrak{g} \rightarrow \mathfrak{h}$ such that
$$
\Phi\left(e^{X}\right)=e^{\phi(X)}
$$
for all $X \in \mathfrak{g}$. The map $\phi$ has following additional properties:
\begin{enumerate}
  \item $\phi\left(A X A^{-1}\right)=\Phi(A) \phi(X) \Phi(A)^{-1}$, for all $X \in \mathfrak{g}, A \in G$.
  \item $\phi([X, Y])=[\phi(X), \phi(Y)]$, for all $X, Y \in \mathfrak{g}$.
  \item $\phi(X)=\left.\frac{d}{d t} \Phi\left(e^{t X}\right)\right|_{t=0}$, for all $X \in \mathfrak{g}$.
\end{enumerate}
\end{THEO}

\begin{PROP}{LIE-H15-03-61}{Kompositionalität von induzierten Lie Algebra Homomorphismen}
Suppose that $G, H$, and $K$ are matrix Lie groups and $\Phi: H \rightarrow K$ and $\Psi: G \rightarrow H$ are Lie group homomorphisms. Let $\Lambda: G \rightarrow K$ be the composition of $\Phi$ and $\Psi$ and let $\phi, \psi$, and $\lambda$ be the Lie algebra maps associated to $\Phi, \Psi$, and $\Lambda$, respectively. Then we have
$$
\lambda=\phi \circ \psi .
$$
\end{PROP}

\begin{PROP}{LIE-H15-03-63}{Induzierte Lie Algebra von (abgeschlossenen, normalen) Kern eines LG Homomorphismus}
If $\Phi: G \rightarrow H$ is a Lie group homomorphism and $\phi: \mathfrak{g} \rightarrow \mathfrak{h}$ is the associated Lie algebra homomorphism, then the kernel of $\Phi$ is a closed, normal subgroup of $G$ and the Lie algebra of the kernel is given by
$$
\operatorname{Lie}(\operatorname{ker}(\Phi))=\operatorname{ker}(\phi)
$$
\end{PROP}

\begin{PROP}{LIE-H15-03-67}{LG Homomorphismus von MLG in Gruppe der invertierbaren Homomorphismen der Lie Algebra}
Let $G$ be a matrix Lie group, with Lie algebra $\mathfrak{g}$. Let $\mathrm{GL}(\mathfrak{g})$ denote the group of all invertible linear transformations of $\mathfrak{g}$. Then the map $A \rightarrow \mathrm{Ad}_{A}$ is a homomorphism of $G$ into $\mathrm{GL}(\mathfrak{g})$. Furthermore, for each $A \in G, \operatorname{Ad}_{A}$ satisfies $\operatorname{Ad}_{A}([X, Y])=\left[\operatorname{Ad}_{A}(X), \operatorname{Ad}_{A}(Y)\right]$ for all $X, Y \in \mathfrak{g}$.
\end{PROP}

\begin{PROP}{LIE-H15-03-70}{Konsistenz der adjungierten Abbidlungsbegriffe}
Let $G$ be a matrix Lie group, let $\mathfrak{g}$ be its Lie algebra, and let $\mathrm{Ad}: G \rightarrow \mathrm{GL}(\mathfrak{g})$ be as in Proposition 3.33. Let $\mathrm{ad}: \mathfrak{g} \rightarrow \mathrm{gl}(\mathfrak{g})$ be the associated Lie algebra map. Then for all $X, Y \in \mathfrak{g}$
$$
\operatorname{ad}_{X}(Y)=[X, Y] .
$$
\end{PROP}

\begin{PROP}{LIE-H15-03-72}{$e^XYe^{-X}$ Formel für adjungierte Abbildung auf $M_n(\C)$}
For any $X$ in $M_{n}(\mathbb{C})$, let $\operatorname{ad}_{X}: M_{n}(\mathbb{C}) \rightarrow M_{n}(\mathbb{C})$ be given by $\operatorname{ad}_{X} Y=[X, Y]$. Then for any $Y$ in $M_{n}(\mathbb{C})$, we have
$$
e^{X} Y e^{-X}=\operatorname{Ad}_{e^{X}}(Y)=e^{\operatorname{ad}_{X}}(Y),
$$
where
$$
e^{\operatorname{ad}_{X}}(Y)=Y+[X, Y]+\frac{1}{2}[X,[X, Y]]+\cdots .
$$
\end{PROP}

\begin{PROP}{LIE-H15-03-75}{Komplexifizierte Lie Algebra und eindeutige Erweiterung der Lie-Klammer}
Let $\mathfrak{g}$ be a finite-dimensional real Lie algebra and $\mathfrak{g}_{\mathbb{C}}$ its complexification. Then the bracket operation on $\mathfrak{g}$ has a unique extension to $\mathfrak{g}_{\mathbb{C}}$ that makes $\mathfrak{g}_{\mathbb{C}}$ into a complex Lie algebra. The complex Lie algebra $\mathfrak{g}_{\mathbb{C}}$ is called the complexification of the real Lie algebra $\mathfrak{g}$.
\end{PROP}

\begin{PROP}{LIE-H15-03-77}{Charakterisierung von Komplexifizierung von reellen Lie Algebren durch Matrizen der Form $X+iY$}
Suppose that $\mathfrak{g} \subset M_{n}(\mathbb{C})$ is a real Lie algebra and that for all nonzero $X$ in $\mathfrak{g}$, the element iX is not in $\mathfrak{g}$. Then the "abstract" complexification $\mathfrak{g}_{\mathbb{C}}$ of $\mathfrak{g}$ in Definition 3.36 is isomorphic to the set of matrices in $M_{n}(\mathbb{C})$ that can be expressed in the form $X+i Y$ with $X$ and $Y$ in $\mathfrak{g}$.
\end{PROP}

\begin{PROP}{LIE-H15-03-81}{Eindeutige Erweiterung von reellen Lie Algebra Homomorphismen in komplexe LA nach Komplexifizierung (Universelle Eigenschaft der Komp.)}
Let $\mathfrak{g}$ be a real Lie algebra, $\mathfrak{g}_{\mathbb{C}}$ its complexification, and $\mathfrak{h}$ an arbitrary complex Lie algebra. Then every real Lie algebra homomorphism of $\mathfrak{g}$ into $\mathfrak{h}$ extends uniquely to a complex Lie algebra homomorphism of $\mathfrak{g}_{\mathbb{C}}$ into $\mathfrak{h}$.
\end{PROP}

\begin{THEO}{LIE-H15-03-85}{Charakterisierung von Inklusion in MLG und Lie Algebren durch Logarithmus}
For $0<\varepsilon<\log 2$, let $U_{\varepsilon}=\left\{X \in M_{n}(\mathbb{C}) \mid\|X\|<\varepsilon\right\}$ and let $V_{\varepsilon}=\exp \left(U_{\varepsilon}\right)$. Suppose $G \subset \mathrm{GL}(n ; \mathbb{C})$ is a matrix Lie group with Lie algebra $\mathfrak{g}$. Then there exists $\varepsilon \in(0, \log 2)$ such that for all $A \in V_{\varepsilon}, A$ is in $G$ if and only if $\log A$ is in $\mathfrak{g}$.
\end{THEO}

\begin{LEM}{LIE-H15-03-88}{Inklusion Charakterisierung für Lie Algebren durch logarithmisierte Folgen}
Suppose $B_{m}$ are elements of $G$ and that $B_{m} \rightarrow I$. Let $Y_{m}=\log B_{m}$, which is defined for all sufficiently large $m$. Suppose that $Y_{m}$ is nonzero for all $m$ and that $Y_{m} /\left\|Y_{m}\right\| \rightarrow Y \in M_{n}(\mathbb{C})$. Then $Y$ is in $\mathfrak{g}$.
\end{LEM}

\begin{KORO}{LIE-H15-03-90}{Exponential Abbildung homöomorph für hinreichend kleine Umgebungen von $0$ und $I$}
If $G$ is a matrix Lie group with Lie algebra $\mathfrak{g}$, there exists a neighborhood $U$ of 0 in $\mathfrak{g}$ and a neighborhood $V$ of $I$ in $G$ such that the exponential map takes $U$ homeomorphically onto $V$.
\end{KORO}

\begin{KORO}{LIE-H15-03-92}{Lie Gruppe mit k-dim reeller Lie Algebra als glatte eingebettete $k$-Untermannigfaltigkeit von $M_N(\C)$}
Let $G$ be a matrix Lie group with Lie algebra $\mathfrak{g}$ and let $k$ be the dimension of $\mathfrak{g}$ as a real vector space. Then $G$ is a smooth embedded submanifold of $M_{n}(\mathbb{C})$ of dimension $k$ and hence a Lie group.
\end{KORO}

\begin{KORO}{LIE-H15-03-94}{Matrix Lie Gruppe ist WZSH und ZSH}
It follows from the corollary that $G$ is locally path connected: every point in $G$ has a neighborhood $U$ that is homeomorphic to a ball in $\mathbb{R}^{k}$ and hence path connected. It then follows that $G$ is connected (in the usual topological sense) if and only if it is path connected. (See, for example, Proposition 3.4.25 of [Run].)
\end{KORO}

\begin{KORO}{LIE-H15-03-95}{Lie Algebra von Matrix Lie Gruppen ist der Tangentialraum am Einheitselement}
Suppose $G \subset \mathrm{GL}(n ; \mathbb{C})$ is a matrix Lie group with Lie algebra $\mathfrak{g}$. Then a matrix $X$ is in $\mathfrak{g}$ if and only if there exists a smooth curve $\gamma$ in $M_{n}(\mathbb{C})$ with $\gamma(t) \in G$ for all $t$ and such that $\gamma(0)=I$ and $d \gamma /\left.d t\right|_{t=0}=X$. Thus, $\mathfrak{g}$ is the tangent space at the identity to $G$.
\end{KORO}

\begin{KORO}{LIE-H15-03-97}{Darstellung von Elementen von ZSH Matrix Lie Gruppen durch Exponentiale von Lie Algebren Elemente}
If $G$ is a connected matrix Lie group, every element $A$ of $G$ can be written in the form
$$
A=e^{X_{1}} e^{X_{2}} \cdots e^{X_{m}}
$$
for some $X_{1}, X_{2}, \ldots, X_{m}$ in $\mathfrak{g}$.
\end{KORO}

\begin{LEM}{LIE-H15-03-99}{Kontrollierbarkeit von $A(s)A(t)^{-1}-I$}
Suppose $A:[a, b] \rightarrow \mathrm{GI}(n ; \mathbb{C})$ is a continuous map. Then for all $\varepsilon>0$ there exists $\delta>0$ such that if $s, t \in[a, b]$ satisfy $|s-t|<\delta$, then
$$
\left\|A(s) A(t)^{-1}-I\right\|<\varepsilon .
$$
\end{LEM}

\begin{KORO}{LIE-H15-03-101}{Charakterisierung von Gleichheit von MLG Homomorphismen bei ZSH Definitionsbereich via LA Homomorphismen}
Let $G$ and $H$ be matrix Lie groups with Lie algebras $\mathfrak{g}$ and $\mathfrak{h}$, respectively, and assume $G$ is connected. Suppose $\Phi_{1}$ and $\Phi_{2}$ are Lie group homomorphisms of $G$ into $H$ and that $\phi_{1}$ and $\phi_{2}$ be the associated Lie algebra homomorphisms of $\mathfrak{g}$ into $\mathfrak{h}$. Then if $\phi_{1}=\phi_{2}$, we have $\Phi_{1}=\Phi_{2}$.
\end{KORO}

\begin{KORO}{LIE-H15-03-103}{Stetige Homomorphismen zwischen Matrix Lie Gruppen sind glatt}
Every continuous homomorphism between two matrix Lie groups is smooth.
\end{KORO}

\begin{KORO}{LIE-H15-03-105}{ZSH Matrix Lie Gruppen mit kommutativen Lie Algebren sind kommutativ}
If $G$ is a connected matrix Lie group and the Lie algebra $\mathfrak{g}$ of $G$ is commutative, then $G$ is commutative.
\end{KORO}

\begin{KORO}{LIE-H15-03-107}{Identität-Komponenten von Matrix Lie Gruppen in $M_n(\C)$ sind Matrix Lie Gruppen mit übereinstimmenden Lie Algebren}
If $G \subset M_{n}(\mathbb{C})$ is a matrix Lie group, the identity component $G_{0}$ of $G$ is a closed subgroup of $\mathrm{GI}(n ; \mathbb{C})$ and thus a matrix Lie group. Furthermore, the Lie algebra of $G_{0}$ is the same as the Lie algebra of $G$.
\end{KORO}

\end{document}